\documentclass[11pt]{article}

    \usepackage[breakable]{tcolorbox}
    \usepackage{parskip} % Stop auto-indenting (to mimic markdown behaviour)
    

    % Basic figure setup, for now with no caption control since it's done
    % automatically by Pandoc (which extracts ![](path) syntax from Markdown).
    \usepackage{graphicx}
    % Keep aspect ratio if custom image width or height is specified
    \setkeys{Gin}{keepaspectratio}
    % Maintain compatibility with old templates. Remove in nbconvert 6.0
    \let\Oldincludegraphics\includegraphics
    % Ensure that by default, figures have no caption (until we provide a
    % proper Figure object with a Caption API and a way to capture that
    % in the conversion process - todo).
    \usepackage{caption}
    \DeclareCaptionFormat{nocaption}{}
    \captionsetup{format=nocaption,aboveskip=0pt,belowskip=0pt}

    \usepackage{float}
    \floatplacement{figure}{H} % forces figures to be placed at the correct location
    \usepackage{xcolor} % Allow colors to be defined
    \usepackage{enumerate} % Needed for markdown enumerations to work
    \usepackage{geometry} % Used to adjust the document margins
    \usepackage{amsmath} % Equations
    \usepackage{amssymb} % Equations
    \usepackage{textcomp} % defines textquotesingle
    % Hack from http://tex.stackexchange.com/a/47451/13684:
    \AtBeginDocument{%
        \def\PYZsq{\textquotesingle}% Upright quotes in Pygmentized code
    }
    \usepackage{upquote} % Upright quotes for verbatim code
    \usepackage{eurosym} % defines \euro

    \usepackage{iftex}
    \ifPDFTeX
        \usepackage[T1]{fontenc}
        \IfFileExists{alphabeta.sty}{
              \usepackage{alphabeta}
          }{
              \usepackage[mathletters]{ucs}
              \usepackage[utf8x]{inputenc}
          }
    \else
        \usepackage{fontspec}
        \usepackage{unicode-math}
    \fi
    \usepackage[normalem]{ulem} 
    \usepackage{fancyvrb} % verbatim replacement that allows latex
    \usepackage{grffile} % extends the file name processing of package graphics
                         % to support a larger range
    \makeatletter % fix for old versions of grffile with XeLaTeX
    \@ifpackagelater{grffile}{2019/11/01}
    {
      % Do nothing on new versions
    }
    {
      \def\Gread@@xetex#1{%
        \IfFileExists{"\Gin@base".bb}%
        {\Gread@eps{\Gin@base.bb}}%
        {\Gread@@xetex@aux#1}%
      }
    }
    \makeatother
    \usepackage[Export]{adjustbox} % Used to constrain images to a maximum size
    \adjustboxset{max size={0.9\linewidth}{0.9\paperheight}}

    % The hyperref package gives us a pdf with properly built
    % internal navigation ('pdf bookmarks' for the table of contents,
    % internal cross-reference links, web links for URLs, etc.)
    \usepackage{hyperref}
    % The default LaTeX title has an obnoxious amount of whitespace. By default,
    % titling removes some of it. It also provides customization options.
    \usepackage{titling}
    \usepackage{longtable} % longtable support required by pandoc >1.10
    \usepackage{booktabs}  % table support for pandoc > 1.12.2
    \usepackage{array}     % table support for pandoc >= 2.11.3
    \usepackage{calc}      % table minipage width calculation for pandoc >= 2.11.1
    \usepackage[inline]{enumitem} % IRkernel/repr support (it uses the enumerate* environment)
    \usepackage[normalem]{ulem} % ulem is needed to support strikethroughs (\sout)
                                % normalem makes italics be italics, not underlines
    \usepackage{soul}      % strikethrough (\st) support for pandoc >= 3.0.0
    \usepackage{mathrsfs}
    

    
    % Colors for the hyperref package
    \definecolor{urlcolor}{rgb}{0,.145,.698}
    \definecolor{linkcolor}{rgb}{.71,0.21,0.01}
    \definecolor{citecolor}{rgb}{.12,.54,.11}

    % ANSI colors
    \definecolor{ansi-black}{HTML}{3E424D}
    \definecolor{ansi-black-intense}{HTML}{282C36}
    \definecolor{ansi-red}{HTML}{E75C58}
    \definecolor{ansi-red-intense}{HTML}{B22B31}
    \definecolor{ansi-green}{HTML}{00A250}
    \definecolor{ansi-green-intense}{HTML}{007427}
    \definecolor{ansi-yellow}{HTML}{DDB62B}
    \definecolor{ansi-yellow-intense}{HTML}{B27D12}
    \definecolor{ansi-blue}{HTML}{208FFB}
    \definecolor{ansi-blue-intense}{HTML}{0065CA}
    \definecolor{ansi-magenta}{HTML}{D160C4}
    \definecolor{ansi-magenta-intense}{HTML}{A03196}
    \definecolor{ansi-cyan}{HTML}{60C6C8}
    \definecolor{ansi-cyan-intense}{HTML}{258F8F}
    \definecolor{ansi-white}{HTML}{C5C1B4}
    \definecolor{ansi-white-intense}{HTML}{A1A6B2}
    \definecolor{ansi-default-inverse-fg}{HTML}{FFFFFF}
    \definecolor{ansi-default-inverse-bg}{HTML}{000000}

    % common color for the border for error outputs.
    \definecolor{outerrorbackground}{HTML}{FFDFDF}

    % commands and environments needed by pandoc snippets
    % extracted from the output of `pandoc -s`
    \providecommand{\tightlist}{%
      \setlength{\itemsep}{0pt}\setlength{\parskip}{0pt}}
    \DefineVerbatimEnvironment{Highlighting}{Verbatim}{commandchars=\\\{\}}
    % Add ',fontsize=\small' for more characters per line
    \newenvironment{Shaded}{}{}
    \newcommand{\KeywordTok}[1]{\textcolor[rgb]{0.00,0.44,0.13}{\textbf{{#1}}}}
    \newcommand{\DataTypeTok}[1]{\textcolor[rgb]{0.56,0.13,0.00}{{#1}}}
    \newcommand{\DecValTok}[1]{\textcolor[rgb]{0.25,0.63,0.44}{{#1}}}
    \newcommand{\BaseNTok}[1]{\textcolor[rgb]{0.25,0.63,0.44}{{#1}}}
    \newcommand{\FloatTok}[1]{\textcolor[rgb]{0.25,0.63,0.44}{{#1}}}
    \newcommand{\CharTok}[1]{\textcolor[rgb]{0.25,0.44,0.63}{{#1}}}
    \newcommand{\StringTok}[1]{\textcolor[rgb]{0.25,0.44,0.63}{{#1}}}
    \newcommand{\CommentTok}[1]{\textcolor[rgb]{0.38,0.63,0.69}{\textit{{#1}}}}
    \newcommand{\OtherTok}[1]{\textcolor[rgb]{0.00,0.44,0.13}{{#1}}}
    \newcommand{\AlertTok}[1]{\textcolor[rgb]{1.00,0.00,0.00}{\textbf{{#1}}}}
    \newcommand{\FunctionTok}[1]{\textcolor[rgb]{0.02,0.16,0.49}{{#1}}}
    \newcommand{\RegionMarkerTok}[1]{{#1}}
    \newcommand{\ErrorTok}[1]{\textcolor[rgb]{1.00,0.00,0.00}{\textbf{{#1}}}}
    \newcommand{\NormalTok}[1]{{#1}}

    % Additional commands for more recent versions of Pandoc
    \newcommand{\ConstantTok}[1]{\textcolor[rgb]{0.53,0.00,0.00}{{#1}}}
    \newcommand{\SpecialCharTok}[1]{\textcolor[rgb]{0.25,0.44,0.63}{{#1}}}
    \newcommand{\VerbatimStringTok}[1]{\textcolor[rgb]{0.25,0.44,0.63}{{#1}}}
    \newcommand{\SpecialStringTok}[1]{\textcolor[rgb]{0.73,0.40,0.53}{{#1}}}
    \newcommand{\ImportTok}[1]{{#1}}
    \newcommand{\DocumentationTok}[1]{\textcolor[rgb]{0.73,0.13,0.13}{\textit{{#1}}}}
    \newcommand{\AnnotationTok}[1]{\textcolor[rgb]{0.38,0.63,0.69}{\textbf{\textit{{#1}}}}}
    \newcommand{\CommentVarTok}[1]{\textcolor[rgb]{0.38,0.63,0.69}{\textbf{\textit{{#1}}}}}
    \newcommand{\VariableTok}[1]{\textcolor[rgb]{0.10,0.09,0.49}{{#1}}}
    \newcommand{\ControlFlowTok}[1]{\textcolor[rgb]{0.00,0.44,0.13}{\textbf{{#1}}}}
    \newcommand{\OperatorTok}[1]{\textcolor[rgb]{0.40,0.40,0.40}{{#1}}}
    \newcommand{\BuiltInTok}[1]{{#1}}
    \newcommand{\ExtensionTok}[1]{{#1}}
    \newcommand{\PreprocessorTok}[1]{\textcolor[rgb]{0.74,0.48,0.00}{{#1}}}
    \newcommand{\AttributeTok}[1]{\textcolor[rgb]{0.49,0.56,0.16}{{#1}}}
    \newcommand{\InformationTok}[1]{\textcolor[rgb]{0.38,0.63,0.69}{\textbf{\textit{{#1}}}}}
    \newcommand{\WarningTok}[1]{\textcolor[rgb]{0.38,0.63,0.69}{\textbf{\textit{{#1}}}}}
    \makeatletter
    \newsavebox\pandoc@box
    \newcommand*\pandocbounded[1]{%
      \sbox\pandoc@box{#1}%
      % scaling factors for width and height
      \Gscale@div\@tempa\textheight{\dimexpr\ht\pandoc@box+\dp\pandoc@box\relax}%
      \Gscale@div\@tempb\linewidth{\wd\pandoc@box}%
      % select the smaller of both
      \ifdim\@tempb\p@<\@tempa\p@
        \let\@tempa\@tempb
      \fi
      % scaling accordingly (\@tempa < 1)
      \ifdim\@tempa\p@<\p@
        \scalebox{\@tempa}{\usebox\pandoc@box}%
      % scaling not needed, use as it is
      \else
        \usebox{\pandoc@box}%
      \fi
    }
    \makeatother

    % Define a nice break command that doesn't care if a line doesn't already
    % exist.
    \def\br{\hspace*{\fill} \\* }
    % Math Jax compatibility definitions
    \def\gt{>}
    \def\lt{<}
    \let\Oldtex\TeX
    \let\Oldlatex\LaTeX
    \renewcommand{\TeX}{\textrm{\Oldtex}}
    \renewcommand{\LaTeX}{\textrm{\Oldlatex}}
    % Document parameters
    % Document title
    \title{Trabalho final ALC2}
    \author{Luiz Fernando Bossa\thanks{l.f.bossa@posgrad.ufsc.br} \and Prof. Fermín S. V. Bazán}
    \date{Dezembro 2025}
    
    
    
    
    
    
% Pygments definitions
\makeatletter
\def\PY@reset{\let\PY@it=\relax \let\PY@bf=\relax%
    \let\PY@ul=\relax \let\PY@tc=\relax%
    \let\PY@bc=\relax \let\PY@ff=\relax}
\def\PY@tok#1{\csname PY@tok@#1\endcsname}
\def\PY@toks#1+{\ifx\relax#1\empty\else%
    \PY@tok{#1}\expandafter\PY@toks\fi}
\def\PY@do#1{\PY@bc{\PY@tc{\PY@ul{%
    \PY@it{\PY@bf{\PY@ff{#1}}}}}}}
\def\PY#1#2{\PY@reset\PY@toks#1+\relax+\PY@do{#2}}

\@namedef{PY@tok@w}{\def\PY@tc##1{\textcolor[rgb]{0.73,0.73,0.73}{##1}}}
\@namedef{PY@tok@c}{\let\PY@it=\textit\def\PY@tc##1{\textcolor[rgb]{0.24,0.48,0.48}{##1}}}
\@namedef{PY@tok@cp}{\def\PY@tc##1{\textcolor[rgb]{0.61,0.40,0.00}{##1}}}
\@namedef{PY@tok@k}{\let\PY@bf=\textbf\def\PY@tc##1{\textcolor[rgb]{0.00,0.50,0.00}{##1}}}
\@namedef{PY@tok@kp}{\def\PY@tc##1{\textcolor[rgb]{0.00,0.50,0.00}{##1}}}
\@namedef{PY@tok@kt}{\def\PY@tc##1{\textcolor[rgb]{0.69,0.00,0.25}{##1}}}
\@namedef{PY@tok@o}{\def\PY@tc##1{\textcolor[rgb]{0.40,0.40,0.40}{##1}}}
\@namedef{PY@tok@ow}{\let\PY@bf=\textbf\def\PY@tc##1{\textcolor[rgb]{0.67,0.13,1.00}{##1}}}
\@namedef{PY@tok@nb}{\def\PY@tc##1{\textcolor[rgb]{0.00,0.50,0.00}{##1}}}
\@namedef{PY@tok@nf}{\def\PY@tc##1{\textcolor[rgb]{0.00,0.00,1.00}{##1}}}
\@namedef{PY@tok@nc}{\let\PY@bf=\textbf\def\PY@tc##1{\textcolor[rgb]{0.00,0.00,1.00}{##1}}}
\@namedef{PY@tok@nn}{\let\PY@bf=\textbf\def\PY@tc##1{\textcolor[rgb]{0.00,0.00,1.00}{##1}}}
\@namedef{PY@tok@ne}{\let\PY@bf=\textbf\def\PY@tc##1{\textcolor[rgb]{0.80,0.25,0.22}{##1}}}
\@namedef{PY@tok@nv}{\def\PY@tc##1{\textcolor[rgb]{0.10,0.09,0.49}{##1}}}
\@namedef{PY@tok@no}{\def\PY@tc##1{\textcolor[rgb]{0.53,0.00,0.00}{##1}}}
\@namedef{PY@tok@nl}{\def\PY@tc##1{\textcolor[rgb]{0.46,0.46,0.00}{##1}}}
\@namedef{PY@tok@ni}{\let\PY@bf=\textbf\def\PY@tc##1{\textcolor[rgb]{0.44,0.44,0.44}{##1}}}
\@namedef{PY@tok@na}{\def\PY@tc##1{\textcolor[rgb]{0.41,0.47,0.13}{##1}}}
\@namedef{PY@tok@nt}{\let\PY@bf=\textbf\def\PY@tc##1{\textcolor[rgb]{0.00,0.50,0.00}{##1}}}
\@namedef{PY@tok@nd}{\def\PY@tc##1{\textcolor[rgb]{0.67,0.13,1.00}{##1}}}
\@namedef{PY@tok@s}{\def\PY@tc##1{\textcolor[rgb]{0.73,0.13,0.13}{##1}}}
\@namedef{PY@tok@sd}{\let\PY@it=\textit\def\PY@tc##1{\textcolor[rgb]{0.73,0.13,0.13}{##1}}}
\@namedef{PY@tok@si}{\let\PY@bf=\textbf\def\PY@tc##1{\textcolor[rgb]{0.64,0.35,0.47}{##1}}}
\@namedef{PY@tok@se}{\let\PY@bf=\textbf\def\PY@tc##1{\textcolor[rgb]{0.67,0.36,0.12}{##1}}}
\@namedef{PY@tok@sr}{\def\PY@tc##1{\textcolor[rgb]{0.64,0.35,0.47}{##1}}}
\@namedef{PY@tok@ss}{\def\PY@tc##1{\textcolor[rgb]{0.10,0.09,0.49}{##1}}}
\@namedef{PY@tok@sx}{\def\PY@tc##1{\textcolor[rgb]{0.00,0.50,0.00}{##1}}}
\@namedef{PY@tok@m}{\def\PY@tc##1{\textcolor[rgb]{0.40,0.40,0.40}{##1}}}
\@namedef{PY@tok@gh}{\let\PY@bf=\textbf\def\PY@tc##1{\textcolor[rgb]{0.00,0.00,0.50}{##1}}}
\@namedef{PY@tok@gu}{\let\PY@bf=\textbf\def\PY@tc##1{\textcolor[rgb]{0.50,0.00,0.50}{##1}}}
\@namedef{PY@tok@gd}{\def\PY@tc##1{\textcolor[rgb]{0.63,0.00,0.00}{##1}}}
\@namedef{PY@tok@gi}{\def\PY@tc##1{\textcolor[rgb]{0.00,0.52,0.00}{##1}}}
\@namedef{PY@tok@gr}{\def\PY@tc##1{\textcolor[rgb]{0.89,0.00,0.00}{##1}}}
\@namedef{PY@tok@ge}{\let\PY@it=\textit}
\@namedef{PY@tok@gs}{\let\PY@bf=\textbf}
\@namedef{PY@tok@ges}{\let\PY@bf=\textbf\let\PY@it=\textit}
\@namedef{PY@tok@gp}{\let\PY@bf=\textbf\def\PY@tc##1{\textcolor[rgb]{0.00,0.00,0.50}{##1}}}
\@namedef{PY@tok@go}{\def\PY@tc##1{\textcolor[rgb]{0.44,0.44,0.44}{##1}}}
\@namedef{PY@tok@gt}{\def\PY@tc##1{\textcolor[rgb]{0.00,0.27,0.87}{##1}}}
\@namedef{PY@tok@err}{\def\PY@bc##1{{\setlength{\fboxsep}{\string -\fboxrule}\fcolorbox[rgb]{1.00,0.00,0.00}{1,1,1}{\strut ##1}}}}
\@namedef{PY@tok@kc}{\let\PY@bf=\textbf\def\PY@tc##1{\textcolor[rgb]{0.00,0.50,0.00}{##1}}}
\@namedef{PY@tok@kd}{\let\PY@bf=\textbf\def\PY@tc##1{\textcolor[rgb]{0.00,0.50,0.00}{##1}}}
\@namedef{PY@tok@kn}{\let\PY@bf=\textbf\def\PY@tc##1{\textcolor[rgb]{0.00,0.50,0.00}{##1}}}
\@namedef{PY@tok@kr}{\let\PY@bf=\textbf\def\PY@tc##1{\textcolor[rgb]{0.00,0.50,0.00}{##1}}}
\@namedef{PY@tok@bp}{\def\PY@tc##1{\textcolor[rgb]{0.00,0.50,0.00}{##1}}}
\@namedef{PY@tok@fm}{\def\PY@tc##1{\textcolor[rgb]{0.00,0.00,1.00}{##1}}}
\@namedef{PY@tok@vc}{\def\PY@tc##1{\textcolor[rgb]{0.10,0.09,0.49}{##1}}}
\@namedef{PY@tok@vg}{\def\PY@tc##1{\textcolor[rgb]{0.10,0.09,0.49}{##1}}}
\@namedef{PY@tok@vi}{\def\PY@tc##1{\textcolor[rgb]{0.10,0.09,0.49}{##1}}}
\@namedef{PY@tok@vm}{\def\PY@tc##1{\textcolor[rgb]{0.10,0.09,0.49}{##1}}}
\@namedef{PY@tok@sa}{\def\PY@tc##1{\textcolor[rgb]{0.73,0.13,0.13}{##1}}}
\@namedef{PY@tok@sb}{\def\PY@tc##1{\textcolor[rgb]{0.73,0.13,0.13}{##1}}}
\@namedef{PY@tok@sc}{\def\PY@tc##1{\textcolor[rgb]{0.73,0.13,0.13}{##1}}}
\@namedef{PY@tok@dl}{\def\PY@tc##1{\textcolor[rgb]{0.73,0.13,0.13}{##1}}}
\@namedef{PY@tok@s2}{\def\PY@tc##1{\textcolor[rgb]{0.73,0.13,0.13}{##1}}}
\@namedef{PY@tok@sh}{\def\PY@tc##1{\textcolor[rgb]{0.73,0.13,0.13}{##1}}}
\@namedef{PY@tok@s1}{\def\PY@tc##1{\textcolor[rgb]{0.73,0.13,0.13}{##1}}}
\@namedef{PY@tok@mb}{\def\PY@tc##1{\textcolor[rgb]{0.40,0.40,0.40}{##1}}}
\@namedef{PY@tok@mf}{\def\PY@tc##1{\textcolor[rgb]{0.40,0.40,0.40}{##1}}}
\@namedef{PY@tok@mh}{\def\PY@tc##1{\textcolor[rgb]{0.40,0.40,0.40}{##1}}}
\@namedef{PY@tok@mi}{\def\PY@tc##1{\textcolor[rgb]{0.40,0.40,0.40}{##1}}}
\@namedef{PY@tok@il}{\def\PY@tc##1{\textcolor[rgb]{0.40,0.40,0.40}{##1}}}
\@namedef{PY@tok@mo}{\def\PY@tc##1{\textcolor[rgb]{0.40,0.40,0.40}{##1}}}
\@namedef{PY@tok@ch}{\let\PY@it=\textit\def\PY@tc##1{\textcolor[rgb]{0.24,0.48,0.48}{##1}}}
\@namedef{PY@tok@cm}{\let\PY@it=\textit\def\PY@tc##1{\textcolor[rgb]{0.24,0.48,0.48}{##1}}}
\@namedef{PY@tok@cpf}{\let\PY@it=\textit\def\PY@tc##1{\textcolor[rgb]{0.24,0.48,0.48}{##1}}}
\@namedef{PY@tok@c1}{\let\PY@it=\textit\def\PY@tc##1{\textcolor[rgb]{0.24,0.48,0.48}{##1}}}
\@namedef{PY@tok@cs}{\let\PY@it=\textit\def\PY@tc##1{\textcolor[rgb]{0.24,0.48,0.48}{##1}}}

\def\PYZbs{\char`\\}
\def\PYZus{\char`\_}
\def\PYZob{\char`\{}
\def\PYZcb{\char`\}}
\def\PYZca{\char`\^}
\def\PYZam{\char`\&}
\def\PYZlt{\char`\<}
\def\PYZgt{\char`\>}
\def\PYZsh{\char`\#}
\def\PYZpc{\char`\%}
\def\PYZdl{\char`\$}
\def\PYZhy{\char`\-}
\def\PYZsq{\char`\'}
\def\PYZdq{\char`\"}
\def\PYZti{\char`\~}
% for compatibility with earlier versions
\def\PYZat{@}
\def\PYZlb{[}
\def\PYZrb{]}
\makeatother


    % For linebreaks inside Verbatim environment from package fancyvrb.
    \makeatletter
        \newbox\Wrappedcontinuationbox
        \newbox\Wrappedvisiblespacebox
        \newcommand*\Wrappedvisiblespace {\textcolor{red}{\textvisiblespace}}
        \newcommand*\Wrappedcontinuationsymbol {\textcolor{red}{\llap{\tiny$\m@th\hookrightarrow$}}}
        \newcommand*\Wrappedcontinuationindent {3ex }
        \newcommand*\Wrappedafterbreak {\kern\Wrappedcontinuationindent\copy\Wrappedcontinuationbox}
        % Take advantage of the already applied Pygments mark-up to insert
        % potential linebreaks for TeX processing.
        %        {, <, #, %, $, ' and ": go to next line.
        %        _, }, ^, &, >, - and ~: stay at end of broken line.
        % Use of \textquotesingle for straight quote.
        \newcommand*\Wrappedbreaksatspecials {%
            \def\PYGZus{\discretionary{\char`\_}{\Wrappedafterbreak}{\char`\_}}%
            \def\PYGZob{\discretionary{}{\Wrappedafterbreak\char`\{}{\char`\{}}%
            \def\PYGZcb{\discretionary{\char`\}}{\Wrappedafterbreak}{\char`\}}}%
            \def\PYGZca{\discretionary{\char`\^}{\Wrappedafterbreak}{\char`\^}}%
            \def\PYGZam{\discretionary{\char`\&}{\Wrappedafterbreak}{\char`\&}}%
            \def\PYGZlt{\discretionary{}{\Wrappedafterbreak\char`\<}{\char`\<}}%
            \def\PYGZgt{\discretionary{\char`\>}{\Wrappedafterbreak}{\char`\>}}%
            \def\PYGZsh{\discretionary{}{\Wrappedafterbreak\char`\#}{\char`\#}}%
            \def\PYGZpc{\discretionary{}{\Wrappedafterbreak\char`\%}{\char`\%}}%
            \def\PYGZdl{\discretionary{}{\Wrappedafterbreak\char`\$}{\char`\$}}%
            \def\PYGZhy{\discretionary{\char`\-}{\Wrappedafterbreak}{\char`\-}}%
            \def\PYGZsq{\discretionary{}{\Wrappedafterbreak\textquotesingle}{\textquotesingle}}%
            \def\PYGZdq{\discretionary{}{\Wrappedafterbreak\char`\"}{\char`\"}}%
            \def\PYGZti{\discretionary{\char`\~}{\Wrappedafterbreak}{\char`\~}}%
        }
        % Some characters . , ; ? ! / are not pygmentized.
        % This macro makes them "active" and they will insert potential linebreaks
        \newcommand*\Wrappedbreaksatpunct {%
            \lccode`\~`\.\lowercase{\def~}{\discretionary{\hbox{\char`\.}}{\Wrappedafterbreak}{\hbox{\char`\.}}}%
            \lccode`\~`\,\lowercase{\def~}{\discretionary{\hbox{\char`\,}}{\Wrappedafterbreak}{\hbox{\char`\,}}}%
            \lccode`\~`\;\lowercase{\def~}{\discretionary{\hbox{\char`\;}}{\Wrappedafterbreak}{\hbox{\char`\;}}}%
            \lccode`\~`\:\lowercase{\def~}{\discretionary{\hbox{\char`\:}}{\Wrappedafterbreak}{\hbox{\char`\:}}}%
            \lccode`\~`\?\lowercase{\def~}{\discretionary{\hbox{\char`\?}}{\Wrappedafterbreak}{\hbox{\char`\?}}}%
            \lccode`\~`\!\lowercase{\def~}{\discretionary{\hbox{\char`\!}}{\Wrappedafterbreak}{\hbox{\char`\!}}}%
            \lccode`\~`\/\lowercase{\def~}{\discretionary{\hbox{\char`\/}}{\Wrappedafterbreak}{\hbox{\char`\/}}}%
            \catcode`\.\active
            \catcode`\,\active
            \catcode`\;\active
            \catcode`\:\active
            \catcode`\?\active
            \catcode`\!\active
            \catcode`\/\active
            \lccode`\~`\~
        }
    \makeatother

    \let\OriginalVerbatim=\Verbatim
    \makeatletter
    \renewcommand{\Verbatim}[1][1]{%
        %\parskip\z@skip
        \sbox\Wrappedcontinuationbox {\Wrappedcontinuationsymbol}%
        \sbox\Wrappedvisiblespacebox {\FV@SetupFont\Wrappedvisiblespace}%
        \def\FancyVerbFormatLine ##1{\hsize\linewidth
            \vtop{\raggedright\hyphenpenalty\z@\exhyphenpenalty\z@
                \doublehyphendemerits\z@\finalhyphendemerits\z@
                \strut ##1\strut}%
        }%
        % If the linebreak is at a space, the latter will be displayed as visible
        % space at end of first line, and a continuation symbol starts next line.
        % Stretch/shrink are however usually zero for typewriter font.
        \def\FV@Space {%
            \nobreak\hskip\z@ plus\fontdimen3\font minus\fontdimen4\font
            \discretionary{\copy\Wrappedvisiblespacebox}{\Wrappedafterbreak}
            {\kern\fontdimen2\font}%
        }%

        % Allow breaks at special characters using \PYG... macros.
        \Wrappedbreaksatspecials
        % Breaks at punctuation characters . , ; ? ! and / need catcode=\active
        \OriginalVerbatim[#1,codes*=\Wrappedbreaksatpunct]%
    }
    \makeatother

    % Exact colors from NB
    \definecolor{incolor}{HTML}{303F9F}
    \definecolor{outcolor}{HTML}{D84315}
    \definecolor{cellborder}{HTML}{CFCFCF}
    \definecolor{cellbackground}{HTML}{F7F7F7}

    % prompt
    \makeatletter
    \newcommand{\boxspacing}{\kern\kvtcb@left@rule\kern\kvtcb@boxsep}
    \makeatother
    \newcommand{\prompt}[4]{
        {\ttfamily\llap{{\color{#2}[#3]:\hspace{3pt}#4}}\vspace{-\baselineskip}}
    }
    

    
    % Prevent overflowing lines due to hard-to-break entities
    \sloppy
    % Setup hyperref package
    \hypersetup{
      breaklinks=true,  % so long urls are correctly broken across lines
      colorlinks=true,
      urlcolor=urlcolor,
      linkcolor=linkcolor,
      citecolor=citecolor,
      }
    % Slightly bigger margins than the latex defaults
    
    \geometry{verbose,tmargin=1in,bmargin=1in,lmargin=1in,rmargin=1in}
    
    

\begin{document}
    
    \maketitle
    
    

    
    \hypertarget{uxe1lgebra-linear-computacional-2---trabalho-final}{%
\section*{Álgebra Linear Computacional 2 - Trabalho
Final}\label{uxe1lgebra-linear-computacional-2---trabalho-final}}

    \hypertarget{questuxe3o-1}{%
\subsection*{Questão 1}\label{questuxe3o-1}}

    Sem perda de generalidade, podemos tomar \(x, y\) normalizados com
\(\|x\|_2=1\) e \(\|y\|_2=1\).

Se \(x = \begin{pmatrix} \alpha \\ u \end{pmatrix}\) com \(\alpha\)
escalar e \(u \in \mathbb{R}^{n-1}\), temos:

\[Ax = \lambda x \iff \begin{bmatrix} \lambda & v^T \\ 0 & T_{22} \end{bmatrix} \begin{bmatrix} \alpha \\ u \end{bmatrix} = \begin{bmatrix} \lambda \alpha + v^T u \\ T_{22} u \end{bmatrix} = \lambda \begin{bmatrix} \alpha \\ u \end{bmatrix}\]

Isso implica no sistema: \[
\begin{cases}
\lambda \alpha + v^T u = \lambda \alpha \implies v^T u = 0 \\
T_{22} u = \lambda u \iff (T_{22} - \lambda I)u = 0
\end{cases}
\]

Como \(T_{22} - \lambda I\) é inversível (pois
\(\lambda \notin \lambda(T_{22})\)), temos \(u = 0\) e logo
\(x = (\alpha, 0) \implies x = e_1\) (considerando \(\alpha=1\)).

Seja \(y = \begin{pmatrix} \beta \\ w \end{pmatrix}\), \(\beta\) escalar
e \(w\in \mathbb{R}^{n-1}\). Da igualdade \(y^* A = \lambda y^*\)

\[(\beta^* \ w^*) \begin{bmatrix} \lambda & v^T \\ 0 & T_{22} \end{bmatrix} = \lambda (\beta^* \ w^*)\]

segue o sistema \[
\begin{cases}
\lambda \beta^* = \lambda \beta^* \\
\beta^* v^T + w^* T_{22} = \lambda w^*
\end{cases}
\]

\(\iff w^*(T_{22} - \lambda I) = -\beta^* v^T \implies w^* = -\beta^* v^T (T_{22} - \lambda I)^{-1}\)

Escolhendo \(\beta = 1\), temos
\(y = \begin{pmatrix} 1 \\ w \end{pmatrix}\).

Veja que \(\|y\|_2^2 = 1^2 + \|w\|_2^2\).

Normalizando \(y\), temos \(\hat{y} = \frac{y}{\|y\|}\).

Agora:

\[\tag{$\dagger$}
s(\lambda) = |\hat{y}^* x| = \frac{|y^* x|}{\|y\|} = \frac{1}{\sqrt{1 + \|w\|_2^2}} = \frac{1}{\sqrt{1 + \|(T_{22} - \lambda I)^{-1} v\|_2^2}}\]

    Conforme visto em sala, para o caso em questão,
\[ \sigma = \text{sep}(\lambda, T_{22}) = \sigma_{\min}(T_{22}-\lambda I)\]

Mas da desigualdade da norma, temos
\[ \|(T_{22} - \lambda I)^{-1} v\|_2^2 \le  \|(T_{22} - \lambda I)^{-1}\|_2^2\| v\|_2^2  = \frac1{\sigma^2_{\min}(T_{22}-\lambda I)}\|v\|_2^2 = \frac{\|v\|^2_2}{\sigma^2}\]

    Assim, substituindo essa desigualdade em \((\dagger)\) temos

\[\frac{1}{\sqrt{1 + \|(T_{22} - \lambda I)^{-1} v\|_2^2}} \le \frac{1}{\sqrt{1 + \frac{\|v\|^2_2}{\sigma^2}}} = \frac{1}{\sqrt{\frac{\sigma^2 + \|v\|^2_2}{\sigma^2}}} =  \frac{\sigma}{\sqrt{\sigma^2 + \|v\|^2_2}} \]

    \hypertarget{questuxe3o-2}{%
\subsection*{Questão 2}\label{questuxe3o-2}}

    Como \(\|q\| = 1\), temos \(q^*q =1\) e podemos escrever
\(\lambda = q*\lambda q\).

Logo
\[\tag{1}\rho - \lambda = q^*\lambda q - q^*Aq = q^*(\lambda I - A)q\]

Como \(v\) é autovetor associado a \(\lambda\), podemos escrever
\((\lambda I - A)v=0\). Somando esse zero em (1) temos

\[\rho - \lambda = q^*(\lambda I - A)(q-v) \tag{2}\]

Tomando o módulo e usando Cauchy-Schwartz e a desigualdade de norma de
operador, temos
\[|\rho-\lambda|\le \|q\|\|(\lambda I - A)(q-v)\| \le 1\cdot \|\lambda I - A\|\|q-v\| \tag{3}\]

Agora, para \(\|\lambda I - A\|\) temos a desigualdade triangular
\[\|\lambda I - A\| \le \|\lambda I\| + \|A\|\]

Como \(\lambda\) é autovalor de \(A\), então \(|\lambda|\le \|A\|\).
Logo \[\|\lambda I - A\| \le \|A\| + \|A\| = 2\|A\| \tag{4}\]

Colando essa desigualdae (4) em (3), segue o resultado.

    \hypertarget{questuxe3o-3}{%
\subsection*{Questão 3}\label{questuxe3o-3}}

    Como os \(\mathsf{v}_j\) são ortogonais, então a norma de \(x_k\) ao
quadrado é a soma dos quadrados das coordenadas:

\[\|x_k\|^2 = \sum_{j=1}^k \left(\frac{\mathsf{u}_j^\top b}{\sigma_j}\right)^2\]

Assim, para \(k+1\) temos

\[\|x_{k+1}\|^2_2 = \sum_{j=1}^{k+1} \left(\frac{\mathsf{u}_j^\top b}{\sigma_j}\right)^2 = \sum_{j=1}^{k} \left(\frac{\mathsf{u}_j^\top b}{\sigma_j}\right)^2 + \left(\frac{\mathsf{u}_{k+1}^\top b}{\sigma_{k+1}}\right)^2  = \|x_k\|^2_2 + \left(\frac{\mathsf{u}_{k+1}^\top b}{\sigma_{k+1}}\right)^2 \ge \|x_k\|^2_2\]

Agora note que o espaço coluna \(\mathcal{S}_k = R(V_k)\) forma uma
sequência encaixada de subespaços, com
\[\mathcal{S}_k\subseteq \mathcal{S}_{k+1}.\]

Para qualquer função positiva, vale que se \(A\subseteq B\), então o
mínimo sobre \(A\) é maior ou igual ao mínimo sobre \(B\). Em
particular, vale para a norma.

Assim,
\[\min_{x\in \mathcal{S}_{k}} \|Ax-b\| \ge \min_{x\in \mathcal{S}_{k+1}} \|Ax-b\|\]
ou seja, \[\|r_{k}\|_2 \ge \|r_{k+1}\|_2\]

    Seja \[A = \sum_{j=1}^n \sigma_ju_jv_j^\top\] a decomposição SVD de
\(A\). Defina \[A^{[k]} = \sum_{j=1}^k \sigma_ju_jv_j^\top \] a
aproximação de posto \(k\) para \(A\). A teoria da SVD garante que
\[\min_{\text{rank}(B)\le k} \|A-B\|_F = \sigma_{k+1}\] e esse valor
mínimo é atingido em \(A^{[k]}\).

Como \(A_k\) tem posto \(k\), então em particular o valor da diferença
da norma de \(A\) e \(A_k\) não pode ser menor do que o \(\sigma_{k+1}\)
acima, logo \[\gamma_k = \|A-A_k\| \ge \sigma_{k+1} \]

Agora note que \(P_k\) é uma matriz de projeção, que projeta em
\(R(V_k) = \mathcal{S}_k\)

Note que \[A-A_k = A-AP_k = A(I-P_k) =: AQ_k\]

com \(Q_k\) sendo a projeção em \(\mathcal{S}_k^\perp\). Como
\(\mathcal{S}_k\subseteq \mathcal{S}_{k+1}\),
\(\mathcal{S}_{k+1}^\perp \subseteq \mathcal{S}_{k}^\perp\).

Assim,

\[\gamma_{k} = \|A-A_{k} \| = \|AQ_{k}\| = \sup^{\|x\|=1}_{x\in\mathcal{S}_{k}^\perp}{\|Ax\|} \ge \sup^{\|x\|=1}_{x\in\mathcal{S}_{k+1}^\perp}{\|Ax\|} = \gamma_{k+1}\]

    \hypertarget{questuxe3o-4}{%
\subsection*{Questão 4}\label{questuxe3o-4}}

Reconstruindo a imagem borrada utilizando a linguagem de programação \texttt{julia}. 

    \begin{tcolorbox}[breakable, size=fbox, boxrule=1pt, pad at break*=1mm,colback=cellbackground, colframe=cellborder]
\prompt{In}{incolor}{ }{\boxspacing}
\begin{Verbatim}[commandchars=\\\{\}]
\PY{k}{using}\PY{+w}{ }\PY{n}{LinearAlgebra}
\PY{k}{using}\PY{+w}{ }\PY{n}{SparseArrays}
\PY{k}{using}\PY{+w}{ }\PY{n}{Random}
\PY{k}{using}\PY{+w}{ }\PY{n}{Plots}
\end{Verbatim}
\end{tcolorbox}

    Instalando pacotes

    \begin{tcolorbox}[breakable, size=fbox, boxrule=1pt, pad at break*=1mm,colback=cellbackground, colframe=cellborder]
\prompt{In}{incolor}{ }{\boxspacing}
\begin{Verbatim}[commandchars=\\\{\}]
\PY{k}{using}\PY{+w}{ }\PY{n}{Pkg}\PY{p}{;}\PY{+w}{ }\PY{n}{Pkg}\PY{o}{.}\PY{n}{add}\PY{p}{(}\PY{p}{[}\PY{l+s}{\PYZdq{}}\PY{l+s}{Images}\PY{l+s}{\PYZdq{}}\PY{p}{,}\PY{+w}{ }\PY{l+s}{\PYZdq{}}\PY{l+s}{HTTP}\PY{l+s}{\PYZdq{}}\PY{p}{]}\PY{p}{,}\PY{+w}{ }\PY{n}{io}\PY{o}{=}\PY{n+nb}{devnull}\PY{p}{)}
\end{Verbatim}
\end{tcolorbox}

    Baixando a imagem

    \begin{tcolorbox}[breakable, size=fbox, boxrule=1pt, pad at break*=1mm,colback=cellbackground, colframe=cellborder]
\prompt{In}{incolor}{ }{\boxspacing}
\begin{Verbatim}[commandchars=\\\{\}]
\PY{k}{using}\PY{+w}{ }\PY{n}{HTTP}
\PY{k}{using}\PY{+w}{ }\PY{n}{Images}

\PY{n}{url}\PY{+w}{ }\PY{o}{=}\PY{+w}{ }\PY{l+s}{\PYZdq{}}\PY{l+s}{https://raw.githubusercontent.com/LFBossa/Disciplinas\PYZhy{}Doutorado/refs/heads/main/AlgebraLinear2/lenaG.png}\PY{l+s}{\PYZdq{}}
\PY{n}{filename}\PY{+w}{ }\PY{o}{=}\PY{+w}{ }\PY{l+s}{\PYZdq{}}\PY{l+s}{lenaG.png}\PY{l+s}{\PYZdq{}}\PY{+w}{ }\PY{c}{\PYZsh{} Use the correct filename}
\PY{n}{HTTP}\PY{o}{.}\PY{n}{download}\PY{p}{(}\PY{n}{url}\PY{p}{,}\PY{+w}{ }\PY{n}{filename}\PY{p}{)}\PY{+w}{ }\PY{c}{\PYZsh{} Download the file}
\PY{n}{lena\PYZus{}img}\PY{+w}{ }\PY{o}{=}\PY{+w}{ }\PY{n}{load}\PY{p}{(}\PY{n}{filename}\PY{p}{)}\PY{+w}{ }\PY{c}{\PYZsh{} Use ImageIO.load to read the image}
\end{Verbatim}
\end{tcolorbox}

    \begin{Verbatim}[commandchars=\\\{\}]
\textcolor{ansi-cyan-intense}{\textbf{┌ }}\textcolor{ansi-cyan-intense}{\textbf{Info: }}Downloading
\textcolor{ansi-cyan-intense}{\textbf{│ }}  source =
"https://raw.githubusercontent.com/LFBossa/Disciplinas-
Doutorado/refs/heads/main/AlgebraLinear2/lenaG.png"
\textcolor{ansi-cyan-intense}{\textbf{│ }}  dest = "lenaG.png"
\textcolor{ansi-cyan-intense}{\textbf{│ }}  progress = 1.0
\textcolor{ansi-cyan-intense}{\textbf{│ }}  time\_taken = "0.1 s"
\textcolor{ansi-cyan-intense}{\textbf{│ }}  time\_remaining = "0.0 s"
\textcolor{ansi-cyan-intense}{\textbf{│ }}  average\_speed = "331.736 KiB/s"
\textcolor{ansi-cyan-intense}{\textbf{│ }}  downloaded = "31.847 KiB"
\textcolor{ansi-cyan-intense}{\textbf{│ }}  remaining = "0 bytes"
\textcolor{ansi-cyan-intense}{\textbf{└ }}  total = "31.847 KiB"
    \end{Verbatim}
 
            
\prompt{Out}{outcolor}{ }{}
    
    \begin{center}
    \adjustimage{max size={0.9\linewidth}{0.9\paperheight}}{Trabalho_final_ALC2_files/Trabalho_final_ALC2_15_1.png}
    \end{center}
    { \hspace*{\fill} \\}
    

    \begin{tcolorbox}[breakable, size=fbox, boxrule=1pt, pad at break*=1mm,colback=cellbackground, colframe=cellborder]
\prompt{In}{incolor}{ }{\boxspacing}
\begin{Verbatim}[commandchars=\\\{\}]
\PY{c}{\PYZsh{} funções auxiliares}
\PY{n}{N}\PY{p}{,}\PY{+w}{ }\PY{n}{N}\PY{+w}{ }\PY{o}{=}\PY{+w}{ }\PY{n}{size}\PY{p}{(}\PY{n}{lena\PYZus{}img}\PY{p}{)}

\PY{n}{imagem2vetor}\PY{p}{(}\PY{n}{img}\PY{p}{)}\PY{+w}{ }\PY{o}{=}\PY{+w}{ }\PY{n}{convert}\PY{p}{(}\PY{k+kt}{Array}\PY{p}{\PYZob{}}\PY{k+kt}{Float64}\PY{p}{\PYZcb{}}\PY{p}{,}\PY{n}{img}\PY{p}{[}\PY{o}{:}\PY{p}{]}\PY{p}{)}
\PY{n}{vetor2imagem}\PY{p}{(}\PY{n}{vetor}\PY{p}{)}\PY{+w}{ }\PY{o}{=}\PY{+w}{ }\PY{n}{Gray}\PY{o}{.}\PY{p}{(}\PY{n}{reshape}\PY{p}{(}\PY{n}{vetor}\PY{p}{,}\PY{+w}{ }\PY{n}{N}\PY{p}{,}\PY{+w}{ }\PY{n}{N}\PY{p}{)}\PY{p}{)}

\PY{n}{lena\PYZus{}vec}\PY{+w}{ }\PY{o}{=}\PY{+w}{ }\PY{n}{imagem2vetor}\PY{p}{(}\PY{n}{lena\PYZus{}img}\PY{p}{)}\PY{p}{;}
\end{Verbatim}
\end{tcolorbox}

    Função para criar a matriz de blur.

    \begin{tcolorbox}[breakable, size=fbox, boxrule=1pt, pad at break*=1mm,colback=cellbackground, colframe=cellborder]
\prompt{In}{incolor}{ }{\boxspacing}
\begin{Verbatim}[commandchars=\\\{\}]
\PY{l+s}{\PYZdq{}\PYZdq{}\PYZdq{}}
\PY{l+s}{    blur(N, band=3, sigma=0.7)}

\PY{l+s}{Problema de teste: desfoque de imagem digital (digital image deblurring).}

\PY{l+s}{A matriz A é uma matriz simétrica duplamente bloco Toeplitz de dimensão N*N\PYZhy{}por\PYZhy{}N*N}
\PY{l+s}{que modela o desfoque de uma imagem N\PYZhy{}por\PYZhy{}N por uma função de espalhamento de ponto Gaussiana.}
\PY{l+s}{Ela é armazenada no formato de matriz esparsa.}

\PY{l+s}{Em cada bloco Toeplitz, apenas elementos da matriz dentro de uma distância band\PYZhy{}1}
\PY{l+s}{da diagonal são não\PYZhy{}nulos (ou seja, band é a meia\PYZhy{}largura de banda).}
\PY{l+s}{Se band não for especificado, band = 3 é usado.}

\PY{l+s}{O parâmetro sigma controla a largura da função de espalhamento de ponto Gaussiana}
\PY{l+s}{e, portanto, a quantidade de suavização (quanto maior o sigma, mais larga a função}
\PY{l+s}{e mais mal posto o problema). Se sigma não for especificado, sigma = 0.7 é usado.}
\PY{l+s}{\PYZdq{}\PYZdq{}\PYZdq{}}
\PY{k}{function}\PY{+w}{ }\PY{n}{blur}\PY{p}{(}\PY{n}{N}\PY{o}{::}\PY{k+kt}{Int}\PY{p}{,}\PY{+w}{ }\PY{n}{band}\PY{o}{::}\PY{k+kt}{Int}\PY{o}{=}\PY{l+m+mi}{3}\PY{p}{,}\PY{+w}{ }\PY{n}{sigma}\PY{o}{::}\PY{k+kt}{Real}\PY{o}{=}\PY{l+m+mf}{0.7}\PY{p}{)}
\PY{+w}{    }\PY{c}{\PYZsh{} Initialization.}
\PY{+w}{    }\PY{n}{band}\PY{+w}{ }\PY{o}{=}\PY{+w}{ }\PY{n}{min}\PY{p}{(}\PY{n}{band}\PY{p}{,}\PY{+w}{ }\PY{n}{N}\PY{p}{)}

\PY{+w}{    }\PY{c}{\PYZsh{} Construct the matrix as a Kronecker product.}
\PY{+w}{    }\PY{c}{\PYZsh{} z vector contains the Gaussian values for the diagonals.}
\PY{+w}{    }\PY{c}{\PYZsh{} MATLAB: z = [exp(\PYZhy{}((0:band\PYZhy{}1).\PYZca{}2)/(2*sigma\PYZca{}2)),zeros(1,N\PYZhy{}band)];}
\PY{+w}{    }\PY{n}{z\PYZus{}vals}\PY{+w}{ }\PY{o}{=}\PY{+w}{ }\PY{n}{exp}\PY{o}{.}\PY{p}{(}\PY{o}{\PYZhy{}}\PY{p}{(}\PY{p}{(}\PY{l+m+mi}{0}\PY{o}{:}\PY{n}{band}\PY{o}{\PYZhy{}}\PY{l+m+mi}{1}\PY{p}{)}\PY{o}{.\PYZca{}}\PY{l+m+mi}{2}\PY{p}{)}\PY{+w}{ }\PY{o}{./}\PY{+w}{ }\PY{p}{(}\PY{l+m+mi}{2}\PY{o}{*}\PY{n}{sigma}\PY{o}{\PYZca{}}\PY{l+m+mi}{2}\PY{p}{)}\PY{p}{)}

\PY{+w}{    }\PY{c}{\PYZsh{} Construct the 1D Toeplitz matrix using spdiagm}
\PY{+w}{    }\PY{c}{\PYZsh{} Diagonal 0 gets z\PYZus{}vals[1]}
\PY{+w}{    }\PY{c}{\PYZsh{} Diagonals +k and \PYZhy{}k get z\PYZus{}vals[k+1]}
\PY{+w}{    }\PY{n}{diags}\PY{+w}{ }\PY{o}{=}\PY{+w}{ }\PY{k+kt}{Pair}\PY{p}{\PYZob{}}\PY{k+kt}{Int}\PY{p}{,}\PY{+w}{ }\PY{k+kt}{Vector}\PY{p}{\PYZob{}}\PY{k+kt}{Float64}\PY{p}{\PYZcb{}}\PY{p}{\PYZcb{}}\PY{p}{[}\PY{p}{]}
\PY{+w}{    }\PY{n}{push!}\PY{p}{(}\PY{n}{diags}\PY{p}{,}\PY{+w}{ }\PY{l+m+mi}{0}\PY{+w}{ }\PY{o}{=\PYZgt{}}\PY{+w}{ }\PY{n}{fill}\PY{p}{(}\PY{n}{z\PYZus{}vals}\PY{p}{[}\PY{l+m+mi}{1}\PY{p}{]}\PY{p}{,}\PY{+w}{ }\PY{n}{N}\PY{p}{)}\PY{p}{)}

\PY{+w}{    }\PY{k}{for}\PY{+w}{ }\PY{n}{k}\PY{+w}{ }\PY{k}{in}\PY{+w}{ }\PY{l+m+mi}{1}\PY{o}{:}\PY{p}{(}\PY{n}{band}\PY{o}{\PYZhy{}}\PY{l+m+mi}{1}\PY{p}{)}
\PY{+w}{        }\PY{n}{val}\PY{+w}{ }\PY{o}{=}\PY{+w}{ }\PY{n}{z\PYZus{}vals}\PY{p}{[}\PY{n}{k}\PY{o}{+}\PY{l+m+mi}{1}\PY{p}{]}
\PY{+w}{        }\PY{n}{push!}\PY{p}{(}\PY{n}{diags}\PY{p}{,}\PY{+w}{ }\PY{n}{k}\PY{+w}{ }\PY{o}{=\PYZgt{}}\PY{+w}{ }\PY{n}{fill}\PY{p}{(}\PY{n}{val}\PY{p}{,}\PY{+w}{ }\PY{n}{N}\PY{o}{\PYZhy{}}\PY{n}{k}\PY{p}{)}\PY{p}{)}
\PY{+w}{        }\PY{n}{push!}\PY{p}{(}\PY{n}{diags}\PY{p}{,}\PY{+w}{ }\PY{o}{\PYZhy{}}\PY{n}{k}\PY{+w}{ }\PY{o}{=\PYZgt{}}\PY{+w}{ }\PY{n}{fill}\PY{p}{(}\PY{n}{val}\PY{p}{,}\PY{+w}{ }\PY{n}{N}\PY{o}{\PYZhy{}}\PY{n}{k}\PY{p}{)}\PY{p}{)}
\PY{+w}{    }\PY{k}{end}

\PY{+w}{    }\PY{n}{A\PYZus{}1d}\PY{+w}{ }\PY{o}{=}\PY{+w}{ }\PY{n}{spdiagm}\PY{p}{(}\PY{n}{N}\PY{p}{,}\PY{+w}{ }\PY{n}{N}\PY{p}{,}\PY{+w}{ }\PY{n}{diags}\PY{o}{...}\PY{p}{)}

\PY{+w}{    }\PY{c}{\PYZsh{} A = (1/(2*pi*sigma\PYZca{}2))*kron(A,A);}
\PY{+w}{    }\PY{n}{A}\PY{+w}{ }\PY{o}{=}\PY{+w}{ }\PY{p}{(}\PY{l+m+mi}{1}\PY{o}{/}\PY{p}{(}\PY{l+m+mi}{2}\PY{o}{*}\PY{n+nb}{pi}\PY{o}{*}\PY{n}{sigma}\PY{o}{\PYZca{}}\PY{l+m+mi}{2}\PY{p}{)}\PY{p}{)}\PY{+w}{ }\PY{o}{*}\PY{+w}{ }\PY{n}{kron}\PY{p}{(}\PY{n}{A\PYZus{}1d}\PY{p}{,}\PY{+w}{ }\PY{n}{A\PYZus{}1d}\PY{p}{)}


\PY{+w}{    }\PY{k}{return}\PY{+w}{ }\PY{n}{A}
\PY{k}{end}
\end{Verbatim}
\end{tcolorbox}

            \begin{tcolorbox}[breakable, size=fbox, boxrule=.5pt, pad at break*=1mm, opacityfill=0]
\prompt{Out}{outcolor}{ }{\boxspacing}
\begin{Verbatim}[commandchars=\\\{\}]
blur
\end{Verbatim}
\end{tcolorbox}
        
    \begin{tcolorbox}[breakable, size=fbox, boxrule=1pt, pad at break*=1mm,colback=cellbackground, colframe=cellborder]
\prompt{In}{incolor}{ }{\boxspacing}
\begin{Verbatim}[commandchars=\\\{\}]
\PY{k}{function}\PY{+w}{ }\PY{n}{ortogonalizar}\PY{p}{(}\PY{n}{U}\PY{o}{::}\PY{k+kt}{Vector}\PY{p}{\PYZob{}}\PY{o}{\PYZlt{}:}\PY{k+kt}{AbstractVector}\PY{p}{\PYZcb{}}\PY{p}{,}\PY{+w}{ }\PY{n}{v}\PY{o}{::}\PY{k+kt}{AbstractVector}\PY{p}{)}
\PY{+w}{  }\PY{c}{\PYZsh{} Ortogonaliza v diante do conjunto Q, usando Gram\PYZhy{}Schimidt}
\PY{+w}{    }\PY{n}{u}\PY{+w}{ }\PY{o}{=}\PY{+w}{ }\PY{n}{copy}\PY{p}{(}\PY{n}{v}\PY{p}{)}
\PY{+w}{    }\PY{k}{for}\PY{+w}{ }\PY{n}{q}\PY{+w}{ }\PY{k}{in}\PY{+w}{ }\PY{n}{U}
\PY{+w}{        }\PY{n}{u}\PY{+w}{ }\PY{o}{\PYZhy{}=}\PY{+w}{ }\PY{p}{(}\PY{n}{dot}\PY{p}{(}\PY{n}{q}\PY{p}{,}\PY{+w}{ }\PY{n}{u}\PY{p}{)}\PY{p}{)}\PY{+w}{ }\PY{o}{*}\PY{+w}{ }\PY{n}{q}
\PY{+w}{    }\PY{k}{end}
\PY{+w}{    }\PY{n}{nrm}\PY{+w}{ }\PY{o}{=}\PY{+w}{ }\PY{n}{norm}\PY{p}{(}\PY{n}{u}\PY{p}{)}
\PY{+w}{    }\PY{k}{return}\PY{+w}{ }\PY{n}{nrm}\PY{+w}{ }\PY{o}{\PYZgt{}}\PY{+w}{ }\PY{l+m+mf}{1e\PYZhy{}12}\PY{+w}{ }\PY{o}{?}\PY{+w}{ }\PY{n}{u}\PY{+w}{ }\PY{o}{/}\PY{+w}{ }\PY{n}{nrm}\PY{+w}{ }\PY{o}{:}\PY{+w}{ }\PY{n}{zeros}\PY{p}{(}\PY{n}{length}\PY{p}{(}\PY{n}{u}\PY{p}{)}\PY{p}{)}
\PY{k}{end}

\PY{k}{function}\PY{+w}{ }\PY{n}{ortogonalizar}\PY{p}{(}\PY{n}{V}\PY{o}{::}\PY{k+kt}{Dict}\PY{p}{,}\PY{+w}{ }\PY{n}{v}\PY{o}{::}\PY{k+kt}{AbstractVector}\PY{p}{)}
\PY{+w}{  }\PY{c}{\PYZsh{} Ortogonaliza v diante do conjunto Q, usando Gram\PYZhy{}Schimidt}
\PY{+w}{    }\PY{n}{u}\PY{+w}{ }\PY{o}{=}\PY{+w}{ }\PY{n}{copy}\PY{p}{(}\PY{n}{v}\PY{p}{)}
\PY{+w}{    }\PY{n}{n}\PY{+w}{ }\PY{o}{=}\PY{+w}{ }\PY{n}{length}\PY{p}{(}\PY{n}{V}\PY{p}{)}
\PY{+w}{    }\PY{k}{for}\PY{+w}{ }\PY{n}{i}\PY{+w}{ }\PY{k}{in}\PY{+w}{ }\PY{l+m+mi}{1}\PY{o}{:}\PY{n}{n}\PY{o}{\PYZhy{}}\PY{l+m+mi}{1}\PY{+w}{ }\PY{c}{\PYZsh{}  V começa em 0}
\PY{+w}{        }\PY{n}{q}\PY{+w}{ }\PY{o}{=}\PY{+w}{ }\PY{n}{V}\PY{p}{[}\PY{n}{i}\PY{p}{]}
\PY{+w}{        }\PY{n}{u}\PY{+w}{ }\PY{o}{\PYZhy{}=}\PY{+w}{ }\PY{p}{(}\PY{n}{dot}\PY{p}{(}\PY{n}{q}\PY{p}{,}\PY{+w}{ }\PY{n}{u}\PY{p}{)}\PY{p}{)}\PY{+w}{ }\PY{o}{*}\PY{+w}{ }\PY{n}{q}
\PY{+w}{    }\PY{k}{end}
\PY{+w}{    }\PY{n}{nrm}\PY{+w}{ }\PY{o}{=}\PY{+w}{ }\PY{n}{norm}\PY{p}{(}\PY{n}{u}\PY{p}{)}
\PY{+w}{    }\PY{k}{return}\PY{+w}{ }\PY{n}{nrm}\PY{+w}{ }\PY{o}{\PYZgt{}}\PY{+w}{ }\PY{l+m+mf}{1e\PYZhy{}12}\PY{+w}{ }\PY{o}{?}\PY{+w}{ }\PY{n}{u}\PY{+w}{ }\PY{o}{/}\PY{+w}{ }\PY{n}{nrm}\PY{+w}{ }\PY{o}{:}\PY{+w}{ }\PY{n}{zeros}\PY{p}{(}\PY{n}{length}\PY{p}{(}\PY{n}{u}\PY{p}{)}\PY{p}{)}
\PY{k}{end}
\end{Verbatim}
\end{tcolorbox}

            \begin{tcolorbox}[breakable, size=fbox, boxrule=.5pt, pad at break*=1mm, opacityfill=0]
\prompt{Out}{outcolor}{ }{\boxspacing}
\begin{Verbatim}[commandchars=\\\{\}]
ortogonalizar (generic function with 2 methods)
\end{Verbatim}
\end{tcolorbox}
        
    \begin{tcolorbox}[breakable, size=fbox, boxrule=1pt, pad at break*=1mm,colback=cellbackground, colframe=cellborder]
\prompt{In}{incolor}{ }{\boxspacing}
\begin{Verbatim}[commandchars=\\\{\}]
\PY{k}{function}\PY{+w}{ }\PY{n}{bidiaglanczos}\PY{p}{(}\PY{n}{A}\PY{p}{,}\PY{+w}{ }\PY{n}{u1}\PY{o}{::}\PY{k+kt}{Vector}\PY{p}{,}\PY{+w}{ }\PY{n}{k}\PY{o}{::}\PY{k+kt}{Number}\PY{p}{;}\PY{+w}{ }\PY{n}{ortog}\PY{o}{=}\PY{n+nb}{false}\PY{p}{)}
\PY{+w}{  }\PY{c}{\PYZsh{} realiza k passos da bidiagonalização}
\PY{+w}{  }\PY{n}{m}\PY{p}{,}\PY{+w}{ }\PY{n}{n}\PY{+w}{ }\PY{o}{=}\PY{+w}{ }\PY{n}{size}\PY{p}{(}\PY{n}{A}\PY{p}{)}
\PY{+w}{  }\PY{n}{α}\PY{+w}{ }\PY{o}{=}\PY{+w}{ }\PY{k+kt}{Array}\PY{p}{\PYZob{}}\PY{k+kt}{Float64}\PY{p}{\PYZcb{}}\PY{p}{(}\PY{p}{[}\PY{p}{]}\PY{p}{)}\PY{+w}{ }\PY{c}{\PYZsh{} inicializa um array vazio}
\PY{+w}{  }\PY{n}{β}\PY{+w}{ }\PY{o}{=}\PY{+w}{ }\PY{p}{[}\PY{+w}{ }\PY{l+m+mf}{.0}\PY{+w}{ }\PY{p}{]}
\PY{+w}{  }\PY{n}{U}\PY{+w}{ }\PY{o}{=}\PY{+w}{ }\PY{p}{[}\PY{+w}{ }\PY{n}{u1}\PY{+w}{ }\PY{p}{]}
\PY{+w}{  }\PY{n}{V}\PY{+w}{ }\PY{o}{=}\PY{+w}{ }\PY{k+kt}{Dict}\PY{p}{(}\PY{+w}{ }\PY{l+m+mi}{0}\PY{+w}{ }\PY{o}{=\PYZgt{}}\PY{+w}{ }\PY{n}{zeros}\PY{p}{(}\PY{n}{n}\PY{p}{)}\PY{p}{)}\PY{+w}{ }\PY{c}{\PYZsh{} Começando os índices em 0}
\PY{+w}{  }\PY{k}{for}\PY{+w}{ }\PY{n}{j}\PY{+w}{ }\PY{o}{=}\PY{+w}{ }\PY{l+m+mi}{1}\PY{o}{:}\PY{n}{k}
\PY{+w}{    }\PY{n}{qj}\PY{+w}{ }\PY{o}{=}\PY{+w}{ }\PY{n}{A}\PY{o}{\PYZsq{}}\PY{o}{*}\PY{n}{U}\PY{p}{[}\PY{n}{j}\PY{p}{]}\PY{+w}{ }\PY{o}{\PYZhy{}}\PY{+w}{  }\PY{n}{β}\PY{p}{[}\PY{n}{j}\PY{p}{]}\PY{o}{*}\PY{n}{V}\PY{p}{[}\PY{n}{j}\PY{o}{\PYZhy{}}\PY{l+m+mi}{1}\PY{p}{]}
\PY{+w}{    }\PY{n}{aj}\PY{+w}{ }\PY{o}{=}\PY{+w}{ }\PY{n}{norm}\PY{p}{(}\PY{n}{qj}\PY{p}{)}
\PY{+w}{    }\PY{n}{vj}\PY{+w}{ }\PY{o}{=}\PY{+w}{ }\PY{n}{qj}\PY{o}{/}\PY{n}{aj}
\PY{+w}{    }\PY{n}{push!}\PY{p}{(}\PY{n}{α}\PY{p}{,}\PY{+w}{ }\PY{n}{aj}\PY{p}{)}\PY{+w}{  }\PY{c}{\PYZsh{} salva a\PYZus{}j}
\PY{+w}{    }\PY{k}{if}\PY{+w}{ }\PY{n}{ortog}\PY{+w}{ }\PY{o}{==}\PY{+w}{ }\PY{n+nb}{true}
\PY{+w}{      }\PY{n}{vj}\PY{+w}{ }\PY{o}{=}\PY{+w}{ }\PY{n}{ortogonalizar}\PY{p}{(}\PY{n}{V}\PY{p}{,}\PY{n}{vj}\PY{p}{)}
\PY{+w}{    }\PY{k}{end}
\PY{+w}{    }\PY{n}{V}\PY{p}{[}\PY{n}{j}\PY{p}{]}\PY{+w}{ }\PY{o}{=}\PY{+w}{ }\PY{n}{vj}\PY{+w}{     }\PY{c}{\PYZsh{} salva v\PYZus{}j}

\PY{+w}{    }\PY{n}{pj}\PY{+w}{ }\PY{o}{=}\PY{+w}{ }\PY{n}{A}\PY{o}{*}\PY{n}{V}\PY{p}{[}\PY{n}{j}\PY{p}{]}\PY{+w}{ }\PY{o}{\PYZhy{}}\PY{+w}{ }\PY{n}{aj}\PY{o}{*}\PY{n}{U}\PY{p}{[}\PY{n}{j}\PY{p}{]}
\PY{+w}{    }\PY{n}{bj1}\PY{+w}{ }\PY{o}{=}\PY{+w}{ }\PY{n}{norm}\PY{p}{(}\PY{n}{pj}\PY{p}{)}
\PY{+w}{    }\PY{n}{push!}\PY{p}{(}\PY{n}{β}\PY{p}{,}\PY{+w}{ }\PY{n}{bj1}\PY{p}{)}\PY{+w}{ }\PY{c}{\PYZsh{} salva b\PYZus{}\PYZob{}j+1\PYZcb{}}
\PY{+w}{    }\PY{n}{uj1}\PY{+w}{ }\PY{o}{=}\PY{+w}{ }\PY{n}{pj}\PY{o}{/}\PY{n}{bj1}
\PY{+w}{    }\PY{k}{if}\PY{+w}{ }\PY{n}{ortog}\PY{+w}{ }\PY{o}{==}\PY{+w}{ }\PY{n+nb}{true}
\PY{+w}{      }\PY{n}{uj1}\PY{+w}{ }\PY{o}{=}\PY{+w}{ }\PY{n}{ortogonalizar}\PY{p}{(}\PY{n}{U}\PY{p}{,}\PY{n}{uj1}\PY{p}{)}
\PY{+w}{    }\PY{k}{end}
\PY{+w}{    }\PY{n}{push!}\PY{p}{(}\PY{n}{U}\PY{p}{,}\PY{n}{uj1}\PY{p}{)}\PY{+w}{  }\PY{c}{\PYZsh{} salva u\PYZus{}\PYZob{}j+1\PYZcb{}}
\PY{+w}{  }\PY{k}{end}
\PY{+w}{  }\PY{c}{\PYZsh{} retorna vetores alfa e beta e matrizes U e V}
\PY{+w}{  }\PY{k}{return}\PY{+w}{ }\PY{n}{α}\PY{p}{,}\PY{+w}{ }\PY{n}{β}\PY{p}{[}\PY{l+m+mi}{2}\PY{o}{:}\PY{k}{end}\PY{p}{]}\PY{p}{,}\PY{+w}{ }\PY{n}{hcat}\PY{p}{(}\PY{n}{U}\PY{o}{...}\PY{p}{)}\PY{p}{,}\PY{+w}{ }\PY{n}{hcat}\PY{p}{(}\PY{p}{[}\PY{n}{V}\PY{p}{[}\PY{n}{j}\PY{p}{]}\PY{+w}{ }\PY{k}{for}\PY{+w}{ }\PY{n}{j}\PY{o}{=}\PY{l+m+mi}{1}\PY{o}{:}\PY{n}{k}\PY{p}{]}\PY{o}{...}\PY{p}{)}
\PY{k}{end}

\PY{k}{function}\PY{+w}{ }\PY{n}{bidiaglanczos}\PY{p}{(}\PY{n}{A}\PY{p}{,}\PY{+w}{ }\PY{n}{k}\PY{o}{::}\PY{k+kt}{Number}\PY{p}{;}\PY{+w}{ }\PY{n}{ortog}\PY{o}{=}\PY{n+nb}{false}\PY{p}{)}
\PY{+w}{  }\PY{c}{\PYZsh{} realiza k passos da bidiagonalização}
\PY{+w}{  }\PY{n}{m}\PY{p}{,}\PY{+w}{ }\PY{n}{n}\PY{+w}{ }\PY{o}{=}\PY{+w}{ }\PY{n}{size}\PY{p}{(}\PY{n}{A}\PY{p}{)}
\PY{+w}{  }\PY{n}{u1}\PY{+w}{ }\PY{o}{=}\PY{+w}{ }\PY{n}{rand}\PY{p}{(}\PY{n}{m}\PY{p}{)}\PY{+w}{       }\PY{c}{\PYZsh{} vetor inicial aleatório}
\PY{+w}{  }\PY{n}{u1}\PY{+w}{ }\PY{o}{/=}\PY{+w}{ }\PY{n}{norm}\PY{p}{(}\PY{n}{u1}\PY{p}{)}\PY{+w}{     }\PY{c}{\PYZsh{} normaliza qk}
\PY{+w}{  }\PY{k}{return}\PY{+w}{ }\PY{n}{bidiaglanczos}\PY{p}{(}\PY{n}{A}\PY{p}{,}\PY{+w}{ }\PY{n}{u1}\PY{p}{,}\PY{+w}{ }\PY{n}{k}\PY{p}{;}\PY{+w}{ }\PY{n}{ortog}\PY{p}{)}
\PY{k}{end}



\PY{k}{function}\PY{+w}{ }\PY{n}{bidiaglanczos}\PY{p}{(}\PY{n}{A}\PY{p}{,}\PY{+w}{ }\PY{n}{u1}\PY{o}{::}\PY{k+kt}{Vector}\PY{p}{;}\PY{+w}{ }\PY{n}{ortog}\PY{o}{=}\PY{n+nb}{false}\PY{p}{)}
\PY{+w}{  }\PY{n}{m}\PY{p}{,}\PY{+w}{ }\PY{n}{n}\PY{+w}{ }\PY{o}{=}\PY{+w}{ }\PY{n}{size}\PY{p}{(}\PY{n}{A}\PY{p}{)}
\PY{+w}{  }\PY{k}{return}\PY{+w}{ }\PY{n}{bidiaglanczos}\PY{p}{(}\PY{n}{A}\PY{p}{,}\PY{+w}{ }\PY{n}{u1}\PY{p}{,}\PY{+w}{ }\PY{n}{n}\PY{p}{;}\PY{+w}{ }\PY{n}{ortog}\PY{p}{)}
\PY{k}{end}


\PY{k}{function}\PY{+w}{ }\PY{n}{bidiaglanczos}\PY{p}{(}\PY{n}{A}\PY{p}{;}\PY{+w}{ }\PY{n}{ortog}\PY{o}{=}\PY{n+nb}{false}\PY{p}{)}
\PY{+w}{  }\PY{n}{m}\PY{p}{,}\PY{+w}{ }\PY{n}{n}\PY{+w}{ }\PY{o}{=}\PY{+w}{ }\PY{n}{size}\PY{p}{(}\PY{n}{A}\PY{p}{)}
\PY{+w}{  }\PY{k}{return}\PY{+w}{ }\PY{n}{bidiaglanczos}\PY{p}{(}\PY{n}{A}\PY{p}{,}\PY{n}{n}\PY{p}{;}\PY{+w}{ }\PY{n}{ortog}\PY{p}{)}
\PY{k}{end}
\end{Verbatim}
\end{tcolorbox}

            \begin{tcolorbox}[breakable, size=fbox, boxrule=.5pt, pad at break*=1mm, opacityfill=0]
\prompt{Out}{outcolor}{ }{\boxspacing}
\begin{Verbatim}[commandchars=\\\{\}]
bidiaglanczos (generic function with 4 methods)
\end{Verbatim}
\end{tcolorbox}
        
    \hypertarget{reconstruindo-a-imagem}{%
\subsection*{Reconstruindo a imagem}\label{reconstruindo-a-imagem}}

    \begin{tcolorbox}[breakable, size=fbox, boxrule=1pt, pad at break*=1mm,colback=cellbackground, colframe=cellborder]
\prompt{In}{incolor}{ }{\boxspacing}
\begin{Verbatim}[commandchars=\\\{\}]
\PY{k}{function}\PY{+w}{ }\PY{n}{GKB2}\PY{p}{(}\PY{n}{A}\PY{p}{,}\PY{+w}{ }\PY{n}{b}\PY{o}{::}\PY{k+kt}{Vector}\PY{p}{,}\PY{+w}{ }\PY{n}{k}\PY{o}{::}\PY{k+kt}{Number}\PY{+w}{ }\PY{p}{;}\PY{+w}{ }\PY{n}{ortog}\PY{o}{=}\PY{n+nb}{false}\PY{p}{)}
\PY{+w}{  }\PY{c}{\PYZsh{} realiza k passos da bidiagonalização}
\PY{+w}{  }\PY{n}{tol}\PY{+w}{ }\PY{o}{=}\PY{+w}{ }\PY{l+m+mf}{1e\PYZhy{}14}
\PY{+w}{  }\PY{n}{m}\PY{p}{,}\PY{+w}{ }\PY{n}{n}\PY{+w}{ }\PY{o}{=}\PY{+w}{ }\PY{n}{size}\PY{p}{(}\PY{n}{A}\PY{p}{)}

\PY{+w}{  }\PY{c}{\PYZsh{} Inicializando}
\PY{+w}{  }\PY{n}{β1}\PY{+w}{ }\PY{o}{=}\PY{+w}{ }\PY{n}{norm}\PY{p}{(}\PY{n}{b}\PY{p}{)}
\PY{+w}{  }\PY{n}{u1}\PY{+w}{ }\PY{o}{=}\PY{+w}{ }\PY{n}{b}\PY{o}{/}\PY{n}{β1}
\PY{+w}{  }\PY{n}{r1}\PY{+w}{ }\PY{o}{=}\PY{+w}{ }\PY{n}{A}\PY{o}{\PYZsq{}}\PY{n}{u1}
\PY{+w}{  }\PY{n}{α1}\PY{+w}{ }\PY{o}{=}\PY{+w}{ }\PY{n}{norm}\PY{p}{(}\PY{n}{r1}\PY{p}{)}
\PY{+w}{  }\PY{n}{v1}\PY{+w}{ }\PY{o}{=}\PY{+w}{ }\PY{n}{r1}\PY{o}{/}\PY{n}{α1}

\PY{+w}{  }\PY{n}{α}\PY{+w}{ }\PY{o}{=}\PY{+w}{ }\PY{p}{[}\PY{+w}{ }\PY{n}{α1}\PY{+w}{ }\PY{p}{]}
\PY{+w}{  }\PY{n}{β}\PY{+w}{ }\PY{o}{=}\PY{+w}{ }\PY{p}{[}\PY{+w}{ }\PY{n}{β1}\PY{+w}{ }\PY{p}{]}
\PY{+w}{  }\PY{n}{U}\PY{+w}{ }\PY{o}{=}\PY{+w}{ }\PY{p}{[}\PY{+w}{ }\PY{n}{u1}\PY{+w}{ }\PY{p}{]}
\PY{+w}{  }\PY{n}{V}\PY{+w}{ }\PY{o}{=}\PY{+w}{ }\PY{p}{[}\PY{+w}{ }\PY{n}{v1}\PY{+w}{ }\PY{p}{]}
\PY{+w}{  }\PY{k}{for}\PY{+w}{ }\PY{n}{j}\PY{+w}{ }\PY{o}{=}\PY{+w}{ }\PY{l+m+mi}{1}\PY{o}{:}\PY{n}{k}
\PY{+w}{    }\PY{n}{pj1}\PY{+w}{ }\PY{o}{=}\PY{+w}{ }\PY{n}{A}\PY{o}{*}\PY{n}{V}\PY{p}{[}\PY{n}{j}\PY{p}{]}\PY{+w}{ }\PY{o}{\PYZhy{}}\PY{+w}{ }\PY{n}{α}\PY{p}{[}\PY{n}{j}\PY{p}{]}\PY{o}{*}\PY{n}{U}\PY{p}{[}\PY{n}{j}\PY{p}{]}
\PY{+w}{    }\PY{n}{bj1}\PY{+w}{ }\PY{o}{=}\PY{+w}{ }\PY{n}{norm}\PY{p}{(}\PY{n}{pj1}\PY{p}{)}
\PY{+w}{    }\PY{k}{if}\PY{+w}{ }\PY{n}{bj1}\PY{+w}{ }\PY{o}{\PYZlt{}}\PY{+w}{ }\PY{n}{tol}
\PY{+w}{      }\PY{k}{break}
\PY{+w}{    }\PY{k}{else}
\PY{+w}{      }\PY{n}{push!}\PY{p}{(}\PY{n}{β}\PY{p}{,}\PY{+w}{ }\PY{n}{bj1}\PY{p}{)}\PY{+w}{ }\PY{c}{\PYZsh{} atualiza b\PYZus{}\PYZob{}j+1\PYZcb{}}
\PY{+w}{      }\PY{n}{uj1}\PY{+w}{ }\PY{o}{=}\PY{+w}{ }\PY{n}{pj1}\PY{o}{/}\PY{n}{bj1}
\PY{+w}{      }\PY{k}{if}\PY{+w}{ }\PY{n}{ortog}\PY{+w}{ }\PY{o}{==}\PY{+w}{ }\PY{n+nb}{true}
\PY{+w}{        }\PY{n}{uj1}\PY{+w}{ }\PY{o}{=}\PY{+w}{ }\PY{n}{ortogonalizar}\PY{p}{(}\PY{n}{U}\PY{p}{,}\PY{n}{uj1}\PY{p}{)}
\PY{+w}{      }\PY{k}{end}
\PY{+w}{      }\PY{n}{push!}\PY{p}{(}\PY{n}{U}\PY{p}{,}\PY{+w}{ }\PY{n}{uj1}\PY{p}{)}\PY{+w}{ }\PY{c}{\PYZsh{} atualiza u\PYZus{}\PYZob{}j+1\PYZcb{}}
\PY{+w}{      }\PY{n}{rj1}\PY{+w}{ }\PY{o}{=}\PY{+w}{ }\PY{n}{A}\PY{o}{\PYZsq{}}\PY{n}{uj1}\PY{+w}{ }\PY{o}{\PYZhy{}}\PY{+w}{ }\PY{n}{bj1}\PY{o}{*}\PY{n}{V}\PY{p}{[}\PY{n}{j}\PY{p}{]}
\PY{+w}{      }\PY{n}{aj1}\PY{+w}{ }\PY{o}{=}\PY{+w}{ }\PY{n}{norm}\PY{p}{(}\PY{n}{rj1}\PY{p}{)}
\PY{+w}{      }\PY{k}{if}\PY{+w}{ }\PY{n}{aj1}\PY{+w}{ }\PY{o}{\PYZlt{}}\PY{+w}{ }\PY{n}{tol}
\PY{+w}{        }\PY{k}{break}
\PY{+w}{      }\PY{k}{else}
\PY{+w}{        }\PY{n}{push!}\PY{p}{(}\PY{n}{α}\PY{p}{,}\PY{+w}{ }\PY{n}{aj1}\PY{p}{)}\PY{+w}{ }\PY{c}{\PYZsh{} atualiza a\PYZus{}\PYZob{}j+1\PYZcb{}}
\PY{+w}{        }\PY{n}{vj1}\PY{+w}{ }\PY{o}{=}\PY{+w}{ }\PY{n}{rj1}\PY{o}{/}\PY{n}{aj1}
\PY{+w}{        }\PY{k}{if}\PY{+w}{ }\PY{n}{ortog}\PY{+w}{ }\PY{o}{==}\PY{+w}{ }\PY{n+nb}{true}
\PY{+w}{          }\PY{n}{vj}\PY{+w}{ }\PY{o}{=}\PY{+w}{ }\PY{n}{ortogonalizar}\PY{p}{(}\PY{n}{V}\PY{p}{,}\PY{n}{vj1}\PY{p}{)}
\PY{+w}{        }\PY{k}{end}
\PY{+w}{        }\PY{n}{push!}\PY{p}{(}\PY{n}{V}\PY{p}{,}\PY{n}{vj1}\PY{p}{)}\PY{+w}{ }\PY{c}{\PYZsh{} atualiza v\PYZus{}\PYZob{}j+1\PYZcb{}}
\PY{+w}{      }\PY{k}{end}
\PY{+w}{    }\PY{k}{end}

\PY{+w}{  }\PY{k}{end}
\PY{+w}{  }\PY{c}{\PYZsh{} retorna vetores alfa e beta e matrizes U e V}
\PY{+w}{  }\PY{k}{return}\PY{+w}{ }\PY{n}{α}\PY{p}{,}\PY{+w}{ }\PY{n}{β}\PY{p}{,}\PY{+w}{ }\PY{n}{hcat}\PY{p}{(}\PY{n}{U}\PY{o}{...}\PY{p}{)}\PY{p}{,}\PY{+w}{ }\PY{n}{hcat}\PY{p}{(}\PY{n}{V}\PY{o}{...}\PY{p}{)}
\PY{k}{end}
\end{Verbatim}
\end{tcolorbox}

            \begin{tcolorbox}[breakable, size=fbox, boxrule=.5pt, pad at break*=1mm, opacityfill=0]
\prompt{Out}{outcolor}{ }{\boxspacing}
\begin{Verbatim}[commandchars=\\\{\}]
GKB2 (generic function with 1 method)
\end{Verbatim}
\end{tcolorbox}
        
    \begin{tcolorbox}[breakable, size=fbox, boxrule=1pt, pad at break*=1mm,colback=cellbackground, colframe=cellborder]
\prompt{In}{incolor}{ }{\boxspacing}
\begin{Verbatim}[commandchars=\\\{\}]
\PY{k}{function}\PY{+w}{ }\PY{n}{Bidiag}\PY{p}{(}\PY{n}{α}\PY{p}{,}\PY{+w}{ }\PY{n}{β}\PY{p}{)}
\PY{+w}{  }\PY{c}{\PYZsh{} retorna uma bidiagonal k+1 x k com alfa e beta}
\PY{+w}{  }\PY{n}{k}\PY{+w}{ }\PY{o}{=}\PY{+w}{ }\PY{n}{length}\PY{p}{(}\PY{n}{α}\PY{p}{)}
\PY{+w}{  }\PY{k}{return}\PY{+w}{ }\PY{p}{[}\PY{n}{Diagonal}\PY{p}{(}\PY{n}{α}\PY{p}{)}\PY{p}{;}\PY{+w}{ }\PY{n}{zeros}\PY{p}{(}\PY{l+m+mi}{1}\PY{p}{,}\PY{n}{k}\PY{p}{)}\PY{p}{]}\PY{+w}{ }\PY{o}{+}\PY{+w}{ }\PY{p}{[}\PY{n}{zeros}\PY{p}{(}\PY{l+m+mi}{1}\PY{p}{,}\PY{n}{k}\PY{p}{)}\PY{p}{;}\PY{+w}{ }\PY{n}{Diagonal}\PY{p}{(}\PY{n}{β}\PY{p}{)}\PY{p}{]}
\PY{k}{end}

\PY{k}{function}\PY{+w}{ }\PY{n}{Bidiag2}\PY{p}{(}\PY{n}{α}\PY{p}{,}\PY{+w}{ }\PY{n}{β}\PY{p}{)}
\PY{+w}{  }\PY{c}{\PYZsh{} retorna uma bidiagonal com alfa e beta}
\PY{+w}{  }\PY{n}{k}\PY{+w}{ }\PY{o}{=}\PY{+w}{ }\PY{n}{length}\PY{p}{(}\PY{n}{α}\PY{p}{)}
\PY{+w}{  }\PY{n}{B}\PY{+w}{ }\PY{o}{=}\PY{+w}{ }\PY{k+kt}{Matrix}\PY{p}{(}\PY{n}{Diagonal}\PY{p}{(}\PY{n}{α}\PY{p}{)}\PY{p}{)}
\PY{+w}{  }\PY{k}{for}\PY{+w}{ }\PY{n}{j}\PY{o}{=}\PY{l+m+mi}{1}\PY{o}{:}\PY{n}{k}\PY{o}{\PYZhy{}}\PY{l+m+mi}{1}
\PY{+w}{    }\PY{n}{B}\PY{p}{[}\PY{n}{j}\PY{o}{+}\PY{l+m+mi}{1}\PY{p}{,}\PY{n}{j}\PY{p}{]}\PY{+w}{ }\PY{o}{=}\PY{+w}{ }\PY{n}{β}\PY{p}{[}\PY{n}{j}\PY{p}{]}
\PY{+w}{  }\PY{k}{end}
\PY{+w}{  }\PY{k}{return}\PY{+w}{ }\PY{n}{B}
\PY{k}{end}
\end{Verbatim}
\end{tcolorbox}

            \begin{tcolorbox}[breakable, size=fbox, boxrule=.5pt, pad at break*=1mm, opacityfill=0]
\prompt{Out}{outcolor}{ }{\boxspacing}
\begin{Verbatim}[commandchars=\\\{\}]
Bidiag2 (generic function with 1 method)
\end{Verbatim}
\end{tcolorbox}
        
    \hypertarget{rotauxe7uxf5es-de-givens}{%
\subsection*{Rotações de Givens}\label{rotauxe7uxf5es-de-givens}}

Para resolver de maneira eficiente o problema

\[\arg\min_{y\in S_k} \|B_ky - \beta_1e_1\|_2 \]

    \begin{tcolorbox}[breakable, size=fbox, boxrule=1pt, pad at break*=1mm,colback=cellbackground, colframe=cellborder]
\prompt{In}{incolor}{ }{\boxspacing}
\begin{Verbatim}[commandchars=\\\{\}]
\PY{k}{function}\PY{+w}{ }\PY{n}{QRGivens}\PY{p}{(}\PY{n}{B}\PY{p}{)}
\PY{+w}{  }\PY{n}{m}\PY{p}{,}\PY{+w}{ }\PY{n}{n}\PY{+w}{ }\PY{o}{=}\PY{+w}{ }\PY{n}{size}\PY{p}{(}\PY{n}{B}\PY{p}{)}
\PY{+w}{  }\PY{n}{Givs}\PY{+w}{ }\PY{o}{=}\PY{+w}{ }\PY{p}{[}\PY{p}{]}
\PY{+w}{  }\PY{n}{GB}\PY{+w}{  }\PY{o}{=}\PY{+w}{ }\PY{n}{copy}\PY{p}{(}\PY{n}{B}\PY{p}{)}
\PY{+w}{  }\PY{k}{for}\PY{+w}{ }\PY{n}{j}\PY{o}{=}\PY{l+m+mi}{1}\PY{o}{:}\PY{n}{min}\PY{p}{(}\PY{n}{m}\PY{o}{\PYZhy{}}\PY{l+m+mi}{1}\PY{p}{,}\PY{n}{n}\PY{o}{\PYZhy{}}\PY{l+m+mi}{1}\PY{p}{)}
\PY{+w}{    }\PY{n}{G}\PY{p}{,}\PY{+w}{ }\PY{n}{r}\PY{+w}{ }\PY{o}{=}\PY{+w}{ }\PY{n}{givens}\PY{p}{(}\PY{n}{GB}\PY{p}{,}\PY{+w}{ }\PY{n}{j}\PY{p}{,}\PY{+w}{ }\PY{n}{j}\PY{o}{+}\PY{l+m+mi}{1}\PY{p}{,}\PY{+w}{ }\PY{n}{j}\PY{p}{)}
\PY{+w}{    }\PY{n}{push!}\PY{p}{(}\PY{n}{Givs}\PY{p}{,}\PY{+w}{ }\PY{n}{G}\PY{p}{)}
\PY{+w}{    }\PY{n}{GB}\PY{+w}{ }\PY{o}{.=}\PY{+w}{ }\PY{n}{G}\PY{o}{*}\PY{n}{GB}
\PY{+w}{  }\PY{k}{end}
\PY{+w}{  }\PY{n}{Givs}
\PY{k}{end}

\PY{k}{function}\PY{+w}{ }\PY{n}{montarQ}\PY{p}{(}\PY{n}{givens}\PY{p}{,}\PY{+w}{ }\PY{n}{n}\PY{p}{)}
\PY{+w}{  }\PY{n}{k}\PY{+w}{ }\PY{o}{=}\PY{+w}{ }\PY{n}{length}\PY{p}{(}\PY{n}{givens}\PY{p}{)}
\PY{+w}{  }\PY{n}{Q}\PY{+w}{ }\PY{o}{=}\PY{+w}{ }\PY{k+kt}{Matrix}\PY{p}{(}\PY{n}{Diagonal}\PY{p}{(}\PY{n}{ones}\PY{p}{(}\PY{n}{n}\PY{p}{)}\PY{p}{)}\PY{p}{)}
\PY{+w}{  }\PY{k}{for}\PY{+w}{ }\PY{n}{j}\PY{o}{=}\PY{l+m+mi}{1}\PY{o}{:}\PY{n}{k}
\PY{+w}{    }\PY{n}{Q}\PY{+w}{ }\PY{o}{.=}\PY{+w}{ }\PY{n}{givens}\PY{p}{[}\PY{n}{j}\PY{p}{]}\PY{o}{*}\PY{n}{Q}
\PY{+w}{  }\PY{k}{end}
\PY{+w}{  }\PY{n}{Q}
\PY{k}{end}

\PY{k}{function}\PY{+w}{ }\PY{n}{resolverGivens}\PY{p}{(}\PY{n}{A}\PY{p}{,}\PY{n}{b}\PY{p}{)}
\PY{+w}{  }\PY{l+s}{\PYZdq{}\PYZdq{}\PYZdq{}}\PY{l+s}{Resolve o sistema Ax=b com A bidiagonal}\PY{l+s}{\PYZdq{}\PYZdq{}\PYZdq{}}
\PY{+w}{  }\PY{n}{m}\PY{p}{,}\PY{n}{n}\PY{+w}{ }\PY{o}{=}\PY{+w}{ }\PY{n}{size}\PY{p}{(}\PY{n}{A}\PY{p}{)}
\PY{+w}{  }\PY{n}{Q}\PY{+w}{ }\PY{o}{=}\PY{+w}{ }\PY{n}{montarQ}\PY{p}{(}\PY{n}{QRGivens}\PY{p}{(}\PY{n}{A}\PY{p}{)}\PY{p}{,}\PY{+w}{ }\PY{n}{m}\PY{p}{)}
\PY{+w}{  }\PY{n}{R}\PY{+w}{ }\PY{o}{=}\PY{+w}{ }\PY{n}{Q}\PY{o}{\PYZsq{}}\PY{n}{A}
\PY{+w}{  }\PY{k}{return}\PY{+w}{ }\PY{n}{R}\PY{o}{\PYZbs{}}\PY{p}{(}\PY{n}{Q}\PY{o}{\PYZsq{}}\PY{n}{b}\PY{p}{)}
\PY{k}{end}
\end{Verbatim}
\end{tcolorbox}

            \begin{tcolorbox}[breakable, size=fbox, boxrule=.5pt, pad at break*=1mm, opacityfill=0]
\prompt{Out}{outcolor}{ }{\boxspacing}
\begin{Verbatim}[commandchars=\\\{\}]
resolverGivens (generic function with 1 method)
\end{Verbatim}
\end{tcolorbox}
        
    \hypertarget{algoritmo-lsqr}{%
\subsection*{Algoritmo LSQR}\label{algoritmo-lsqr}}

Vamos usar a bidiagonalização de Golub-Kahan para obter incrementalmente
matrizes \(V_k, B_k\) e \(U_k\) de modo que \[AV_k = U_{k+1}B_k\] de
modo que podemos reformular o problema

\[\arg\min_{x\in S_k} \|Ax-b\| = \arg\min_{y\in\mathbb{R}^n}\|AV_ky-b\| = \arg\min_{y\in\mathbb{R}^n} \|B_ky-\beta_1e_1\| \tag{$\star$}\]

Como \(B_k\) é bidiagonal inferior, podemos usar rotações de Givens para
fazer a fatoração \(QR = B_k\) e resolvemos o problema acima resolvendo
o sistema triangular\\
\[Rx = Q^\top\beta_1e_1 \tag{G}\]

Esse problema \((G)\) é resolvido pela função
\texttt{resolverGivens(Bk,\ β1e1)}

    \begin{tcolorbox}[breakable, size=fbox, boxrule=1pt, pad at break*=1mm,colback=cellbackground, colframe=cellborder]
\prompt{In}{incolor}{ }{\boxspacing}
\begin{Verbatim}[commandchars=\\\{\}]
\PY{k}{function}\PY{+w}{ }\PY{n}{ReconstrucaoLSQR}\PY{p}{(}\PY{n}{A}\PY{p}{,}\PY{+w}{ }\PY{n}{b}\PY{+w}{ }\PY{o}{::}\PY{+w}{ }\PY{k+kt}{Vector}\PY{p}{,}\PY{+w}{ }\PY{n}{δ}\PY{+w}{ }\PY{o}{::}\PY{+w}{ }\PY{k+kt}{Number}\PY{p}{,}\PY{+w}{ }\PY{n}{k}\PY{+w}{ }\PY{o}{::}\PY{+w}{ }\PY{k+kt}{Number}\PY{p}{,}\PY{+w}{ }\PY{n}{τ}\PY{+w}{ }\PY{o}{=}\PY{+w}{ }\PY{l+m+mf}{1.1}\PY{p}{;}\PY{+w}{ }\PY{n}{ortog}\PY{o}{=}\PY{n+nb}{false}\PY{p}{)}
\PY{+w}{  }\PY{c}{\PYZsh{} Realiza a reconstrução da imagem usando o algoritmo GKB2}
\PY{+w}{  }\PY{n}{tol}\PY{+w}{ }\PY{o}{=}\PY{+w}{ }\PY{l+m+mf}{1e\PYZhy{}14}
\PY{+w}{  }\PY{n}{m}\PY{p}{,}\PY{+w}{ }\PY{n}{n}\PY{+w}{ }\PY{o}{=}\PY{+w}{ }\PY{n}{size}\PY{p}{(}\PY{n}{A}\PY{p}{)}
\PY{+w}{  }\PY{n}{β1}\PY{+w}{ }\PY{o}{=}\PY{+w}{ }\PY{n}{norm}\PY{p}{(}\PY{n}{b}\PY{p}{)}
\PY{+w}{  }\PY{n}{u1}\PY{+w}{ }\PY{o}{=}\PY{+w}{ }\PY{n}{b}\PY{o}{/}\PY{n}{β1}
\PY{+w}{  }\PY{n}{r1}\PY{+w}{ }\PY{o}{=}\PY{+w}{ }\PY{n}{A}\PY{o}{\PYZsq{}}\PY{n}{u1}
\PY{+w}{  }\PY{n}{α1}\PY{+w}{ }\PY{o}{=}\PY{+w}{ }\PY{n}{norm}\PY{p}{(}\PY{n}{r1}\PY{p}{)}
\PY{+w}{  }\PY{n}{v1}\PY{+w}{ }\PY{o}{=}\PY{+w}{ }\PY{n}{r1}\PY{o}{/}\PY{n}{α1}

\PY{+w}{  }\PY{n}{α}\PY{+w}{ }\PY{o}{=}\PY{+w}{ }\PY{p}{[}\PY{+w}{ }\PY{n}{α1}\PY{+w}{ }\PY{p}{]}
\PY{+w}{  }\PY{n}{β}\PY{+w}{ }\PY{o}{=}\PY{+w}{ }\PY{p}{[}\PY{+w}{ }\PY{n}{β1}\PY{+w}{ }\PY{p}{]}
\PY{+w}{  }\PY{n}{U}\PY{+w}{ }\PY{o}{=}\PY{+w}{ }\PY{p}{[}\PY{+w}{ }\PY{n}{u1}\PY{+w}{ }\PY{p}{]}
\PY{+w}{  }\PY{n}{V}\PY{+w}{ }\PY{o}{=}\PY{+w}{ }\PY{p}{[}\PY{+w}{ }\PY{n}{v1}\PY{+w}{ }\PY{p}{]}

\PY{+w}{  }\PY{n}{Xizes}\PY{+w}{ }\PY{o}{=}\PY{+w}{ }\PY{p}{[}\PY{p}{]}
\PY{+w}{  }\PY{n}{Erros}\PY{+w}{ }\PY{o}{=}\PY{+w}{ }\PY{p}{[}\PY{p}{]}
\PY{+w}{  }\PY{k}{for}\PY{+w}{ }\PY{n}{j}\PY{+w}{ }\PY{o}{=}\PY{+w}{ }\PY{l+m+mi}{1}\PY{o}{:}\PY{n}{k}
\PY{+w}{    }\PY{n}{pj1}\PY{+w}{ }\PY{o}{=}\PY{+w}{ }\PY{n}{A}\PY{o}{*}\PY{n}{V}\PY{p}{[}\PY{n}{j}\PY{p}{]}\PY{+w}{ }\PY{o}{\PYZhy{}}\PY{+w}{ }\PY{n}{α}\PY{p}{[}\PY{n}{j}\PY{p}{]}\PY{o}{*}\PY{n}{U}\PY{p}{[}\PY{n}{j}\PY{p}{]}
\PY{+w}{    }\PY{n}{bj1}\PY{+w}{ }\PY{o}{=}\PY{+w}{ }\PY{n}{norm}\PY{p}{(}\PY{n}{pj1}\PY{p}{)}
\PY{+w}{    }\PY{k}{if}\PY{+w}{ }\PY{n}{bj1}\PY{+w}{ }\PY{o}{\PYZlt{}}\PY{+w}{ }\PY{n}{tol}\PY{+w}{ }\PY{c}{\PYZsh{} b\PYZus{}\PYZob{}j+1\PYZcb{} = 0}
\PY{+w}{      }\PY{k}{break}
\PY{+w}{    }\PY{k}{else}
\PY{+w}{      }\PY{n}{push!}\PY{p}{(}\PY{n}{β}\PY{p}{,}\PY{+w}{ }\PY{n}{bj1}\PY{p}{)}\PY{+w}{ }\PY{c}{\PYZsh{} atualiza b\PYZus{}\PYZob{}j+1\PYZcb{}}
\PY{+w}{      }\PY{n}{uj1}\PY{+w}{ }\PY{o}{=}\PY{+w}{ }\PY{n}{pj1}\PY{o}{/}\PY{n}{bj1}
\PY{+w}{      }\PY{k}{if}\PY{+w}{ }\PY{n}{ortog}\PY{+w}{ }\PY{o}{==}\PY{+w}{ }\PY{n+nb}{true}
\PY{+w}{        }\PY{n}{uj1}\PY{+w}{ }\PY{o}{=}\PY{+w}{ }\PY{n}{ortogonalizar}\PY{p}{(}\PY{n}{U}\PY{p}{,}\PY{n}{uj1}\PY{p}{)}
\PY{+w}{      }\PY{k}{end}
\PY{+w}{      }\PY{n}{push!}\PY{p}{(}\PY{n}{U}\PY{p}{,}\PY{+w}{ }\PY{n}{uj1}\PY{p}{)}\PY{+w}{ }\PY{c}{\PYZsh{} atualiza u\PYZus{}\PYZob{}j+1\PYZcb{}}
\PY{+w}{      }\PY{n}{rj1}\PY{+w}{ }\PY{o}{=}\PY{+w}{ }\PY{n}{A}\PY{o}{\PYZsq{}}\PY{n}{uj1}\PY{+w}{ }\PY{o}{\PYZhy{}}\PY{+w}{ }\PY{n}{bj1}\PY{o}{*}\PY{n}{V}\PY{p}{[}\PY{n}{j}\PY{p}{]}
\PY{+w}{      }\PY{n}{aj1}\PY{+w}{ }\PY{o}{=}\PY{+w}{ }\PY{n}{norm}\PY{p}{(}\PY{n}{rj1}\PY{p}{)}
\PY{+w}{      }\PY{k}{if}\PY{+w}{ }\PY{n}{aj1}\PY{+w}{ }\PY{o}{\PYZlt{}}\PY{+w}{ }\PY{n}{tol}\PY{+w}{ }\PY{c}{\PYZsh{} a\PYZus{}\PYZob{}j+1\PYZcb{} = 0}
\PY{+w}{        }\PY{k}{break}
\PY{+w}{      }\PY{k}{else}
\PY{+w}{        }\PY{n}{push!}\PY{p}{(}\PY{n}{α}\PY{p}{,}\PY{+w}{ }\PY{n}{aj1}\PY{p}{)}\PY{+w}{ }\PY{c}{\PYZsh{} atualiza a\PYZus{}\PYZob{}j+1\PYZcb{}}
\PY{+w}{        }\PY{n}{vj1}\PY{+w}{ }\PY{o}{=}\PY{+w}{ }\PY{n}{rj1}\PY{o}{/}\PY{n}{aj1}
\PY{+w}{        }\PY{k}{if}\PY{+w}{ }\PY{n}{ortog}\PY{+w}{ }\PY{o}{==}\PY{+w}{ }\PY{n+nb}{true}
\PY{+w}{          }\PY{n}{vj1}\PY{+w}{ }\PY{o}{=}\PY{+w}{ }\PY{n}{ortogonalizar}\PY{p}{(}\PY{n}{V}\PY{p}{,}\PY{n}{vj1}\PY{p}{)}
\PY{+w}{        }\PY{k}{end}
\PY{+w}{        }\PY{n}{push!}\PY{p}{(}\PY{n}{V}\PY{p}{,}\PY{n}{vj1}\PY{p}{)}\PY{+w}{ }\PY{c}{\PYZsh{} atualiza v\PYZus{}\PYZob{}j+1\PYZcb{}}
\PY{+w}{      }\PY{k}{end}
\PY{+w}{    }\PY{k}{end}

\PY{+w}{    }\PY{n}{Vk}\PY{+w}{ }\PY{o}{=}\PY{+w}{ }\PY{n}{hcat}\PY{p}{(}\PY{n}{V}\PY{p}{[}\PY{l+m+mi}{1}\PY{o}{:}\PY{n}{j}\PY{p}{]}\PY{o}{...}\PY{p}{)}
\PY{+w}{    }\PY{n}{Bk}\PY{+w}{ }\PY{o}{=}\PY{+w}{ }\PY{n}{Bidiag}\PY{p}{(}\PY{n}{α}\PY{p}{[}\PY{l+m+mi}{1}\PY{o}{:}\PY{n}{j}\PY{p}{]}\PY{p}{,}\PY{+w}{ }\PY{n}{β}\PY{p}{[}\PY{l+m+mi}{2}\PY{o}{:}\PY{n}{j}\PY{o}{+}\PY{l+m+mi}{1}\PY{p}{]}\PY{p}{)}
\PY{+w}{    }\PY{n}{β1e1}\PY{+w}{ }\PY{o}{=}\PY{+w}{ }\PY{n}{zeros}\PY{p}{(}\PY{n}{j}\PY{o}{+}\PY{l+m+mi}{1}\PY{p}{)}
\PY{+w}{    }\PY{n}{β1e1}\PY{p}{[}\PY{l+m+mi}{1}\PY{p}{]}\PY{+w}{ }\PY{o}{=}\PY{+w}{ }\PY{n}{β1}
\PY{+w}{    }\PY{n}{yk}\PY{+w}{ }\PY{o}{=}\PY{+w}{ }\PY{n}{resolverGivens}\PY{p}{(}\PY{n}{Bk}\PY{p}{,}\PY{n}{β1e1}\PY{p}{)}
\PY{+w}{    }\PY{n}{xk}\PY{+w}{ }\PY{o}{=}\PY{+w}{ }\PY{n}{Vk}\PY{o}{*}\PY{n}{yk}
\PY{+w}{    }\PY{n}{erro}\PY{+w}{ }\PY{o}{=}\PY{+w}{ }\PY{n}{norm}\PY{p}{(}\PY{n}{A}\PY{o}{*}\PY{n}{xk}\PY{o}{\PYZhy{}}\PY{n}{b}\PY{p}{)}
\PY{+w}{    }\PY{n}{push!}\PY{p}{(}\PY{n}{Xizes}\PY{p}{,}\PY{+w}{ }\PY{n}{xk}\PY{p}{)}
\PY{+w}{    }\PY{n}{push!}\PY{p}{(}\PY{n}{Erros}\PY{p}{,}\PY{+w}{ }\PY{n}{erro}\PY{p}{)}
\PY{+w}{    }\PY{k}{if}\PY{+w}{ }\PY{n}{erro}\PY{+w}{ }\PY{o}{<=}\PY{+w}{ }\PY{n}{τ}\PY{o}{*}\PY{n}{δ}
\PY{+w}{      }\PY{n+nd}{@info}\PY{+w}{ }\PY{l+s}{\PYZdq{}}\PY{l+s}{Critério de discrepância aceito com k=}\PY{l+s+si}{\PYZdl{}j}\PY{l+s}{\PYZdq{}}
\PY{+w}{      }\PY{k}{return}\PY{+w}{ }\PY{n}{α}\PY{p}{,}\PY{+w}{ }\PY{n}{β}\PY{p}{,}\PY{+w}{ }\PY{n}{hcat}\PY{p}{(}\PY{n}{U}\PY{o}{...}\PY{p}{)}\PY{p}{,}\PY{+w}{ }\PY{n}{hcat}\PY{p}{(}\PY{n}{V}\PY{o}{...}\PY{p}{)}\PY{p}{,}\PY{+w}{ }\PY{n}{Xizes}\PY{p}{,}\PY{+w}{ }\PY{n}{Erros}
\PY{+w}{    }\PY{k}{end}

\PY{+w}{  }\PY{k}{end}
\PY{+w}{  }\PY{c}{\PYZsh{} retorna vetores alfa e beta e matrizes U e V}
\PY{+w}{     }\PY{n+nd}{@info}\PY{+w}{ }\PY{l+s}{\PYZdq{}}\PY{l+s}{Critério de discrepância não foi atingido}\PY{l+s}{\PYZdq{}}
\PY{+w}{  }\PY{k}{return}\PY{+w}{ }\PY{n}{α}\PY{p}{,}\PY{+w}{ }\PY{n}{β}\PY{p}{,}\PY{+w}{ }\PY{n}{hcat}\PY{p}{(}\PY{n}{U}\PY{o}{...}\PY{p}{)}\PY{p}{,}\PY{+w}{ }\PY{n}{hcat}\PY{p}{(}\PY{n}{V}\PY{o}{...}\PY{p}{)}\PY{p}{,}\PY{+w}{ }\PY{n}{Xizes}\PY{p}{,}\PY{+w}{ }\PY{n}{Erros}
\PY{k}{end}
\end{Verbatim}
\end{tcolorbox}

            \begin{tcolorbox}[breakable, size=fbox, boxrule=.5pt, pad at break*=1mm, opacityfill=0]
\prompt{Out}{outcolor}{ }{\boxspacing}
\begin{Verbatim}[commandchars=\\\{\}]
ReconstrucaoLSQR (generic function with 2 methods)
\end{Verbatim}
\end{tcolorbox}
        
    \begin{Shaded}
\begin{Highlighting}[]
\FunctionTok{ReconstrucaoLSQR}\NormalTok{(A, b, δ, k, τ }\OperatorTok{=} \FloatTok{1.1}\NormalTok{; ortog}\OperatorTok{=}\ConstantTok{false}\NormalTok{)}
\end{Highlighting}
\end{Shaded}

Função que implementa o algoritmo LSQR para reconstrução de imagem,
utilizando a bidiagonalização de Golub-Kahan. Busca uma solução
aproximada \(x_k\) para o sistema linear \(Ax = b\), minimizando
\(||Ax - b||_2\), sujeita a um critério de discrepância, onde \(b\) é a
imagem borrada e \(A\) é a matriz de desfoque.

\textbf{Argumentos}

\begin{itemize}
\tightlist
\item
  \texttt{A}: Matriz de desfoque esparsa (Sparse Matrix), geralmente
  gerada pela função \texttt{blur}.
\item
  \texttt{b}: Vetor da imagem borrada e com ruído.
\item
  \texttt{δ}: Nível de ruído conhecido em \texttt{b}, usado no critério
  de discrepância (\(\|Ax_k - b\| \le \tau\delta\)).
\item
  \texttt{k}: Número máximo de iterações para a bidiagonalização.
\item
  \texttt{τ}: Parâmetro de tolerância para o critério de discrepância
  (padrão: 1.1).
\item
  \texttt{ortog}: Booleano que indica se deve ser aplicada
  ortogonalização explícita (padrão: \texttt{false}).
\end{itemize}

\textbf{Retorna}

\begin{itemize}
\tightlist
\item
  \texttt{α}: Vetor dos elementos diagonais da matriz bidiagonal
  \(B_k\).
\item
  \texttt{β}: Vetor dos elementos subdiagonais da matriz bidiagonal
  \(B_k\).
\item
  \texttt{U}: Matriz ortogonal \(U_k\).
\item
  \texttt{V}: Matriz ortogonal \(V_k\).
\item
  \texttt{Xizes}: Lista de soluções aproximadas \(x_k\) em cada
  iteração.
\item
  \texttt{Erros}: Lista dos resíduos \(\|Ax_k - b\|\) em cada iteração.
\end{itemize}

O algoritmo para quando o critério de discrepância é atendido ou o
número máximo de iterações \texttt{k} é atingido.

    \hypertarget{testes}{%
\subsection*{Testes}\label{testes}}

    \hypertarget{erro-5-tau1.1-com-ortogonalizauxe7uxe3o}{%
\subsubsection*{\texorpdfstring{Erro 5\%, \(\tau=1.1\), com
ortogonalização}{Erro 5\%, \textbackslash tau=1.1, com ortogonalização}}\label{erro-5-tau1.1-com-ortogonalizauxe7uxe3o}}

    \begin{tcolorbox}[breakable, size=fbox, boxrule=1pt, pad at break*=1mm,colback=cellbackground, colframe=cellborder]
\prompt{In}{incolor}{ }{\boxspacing}
\begin{Verbatim}[commandchars=\\\{\}]
\PY{n}{A}\PY{+w}{ }\PY{o}{=}\PY{+w}{  }\PY{n}{blur}\PY{p}{(}\PY{n}{N}\PY{p}{,}\PY{l+m+mi}{5}\PY{p}{,}\PY{l+m+mf}{3.5}\PY{p}{)}\PY{p}{;}
\PY{n}{lena\PYZus{}vec\PYZus{}b}\PY{+w}{ }\PY{o}{=}\PY{+w}{ }\PY{n}{A}\PY{o}{*}\PY{n}{lena\PYZus{}vec}\PY{p}{;}

\PY{n}{erro}\PY{+w}{ }\PY{o}{=}\PY{+w}{ }\PY{n}{randn}\PY{p}{(}\PY{n}{N}\PY{o}{\PYZca{}}\PY{l+m+mi}{2}\PY{p}{)}
\PY{n}{erro}\PY{+w}{ }\PY{o}{=}\PY{+w}{ }\PY{n}{erro}\PY{o}{/}\PY{n}{norm}\PY{p}{(}\PY{n}{erro}\PY{p}{)}

\PY{n}{e}\PY{+w}{ }\PY{o}{=}\PY{+w}{ }\PY{l+m+mf}{0.05}
\PY{n}{τ}\PY{+w}{ }\PY{o}{=}\PY{+w}{ }\PY{l+m+mf}{1.1}

\PY{n}{lena\PYZus{}vec\PYZus{}bn}\PY{+w}{ }\PY{o}{=}\PY{+w}{ }\PY{n}{lena\PYZus{}vec\PYZus{}b}\PY{+w}{ }\PY{o}{+}\PY{+w}{ }\PY{n}{e}\PY{o}{*}\PY{n}{erro}\PY{o}{*}\PY{n}{norm}\PY{p}{(}\PY{n}{lena\PYZus{}vec\PYZus{}b}\PY{p}{)}
\PY{n}{lena\PYZus{}img\PYZus{}bn}\PY{+w}{ }\PY{o}{=}\PY{+w}{ }\PY{n}{vetor2imagem}\PY{p}{(}\PY{n}{lena\PYZus{}vec\PYZus{}bn}\PY{p}{)}
\end{Verbatim}
\end{tcolorbox}
 
            
\prompt{Out}{outcolor}{ }{}
    
    \begin{center}
    \adjustimage{max size={0.9\linewidth}{0.9\paperheight}}{Trabalho_final_ALC2_files/Trabalho_final_ALC2_31_0.png}
    \end{center}
    { \hspace*{\fill} \\}
    

    Critério de discrepância.

    \begin{tcolorbox}[breakable, size=fbox, boxrule=1pt, pad at break*=1mm,colback=cellbackground, colframe=cellborder]
\prompt{In}{incolor}{ }{\boxspacing}
\begin{Verbatim}[commandchars=\\\{\}]
\PY{n}{δ}\PY{+w}{ }\PY{o}{=}\PY{+w}{ }\PY{n}{norm}\PY{p}{(}\PY{n}{lena\PYZus{}vec\PYZus{}b}\PY{+w}{ }\PY{o}{\PYZhy{}}\PY{+w}{ }\PY{n}{lena\PYZus{}vec\PYZus{}bn}\PY{p}{)}
\end{Verbatim}
\end{tcolorbox}

            \begin{tcolorbox}[breakable, size=fbox, boxrule=.5pt, pad at break*=1mm, opacityfill=0]
\prompt{Out}{outcolor}{ }{\boxspacing}
\begin{Verbatim}[commandchars=\\\{\}]
4.21554656170367
\end{Verbatim}
\end{tcolorbox}
        
    \begin{tcolorbox}[breakable, size=fbox, boxrule=1pt, pad at break*=1mm,colback=cellbackground, colframe=cellborder]
\prompt{In}{incolor}{ }{\boxspacing}
\begin{Verbatim}[commandchars=\\\{\}]
\PY{n+nd}{@time}\PY{+w}{ }\PY{n}{α}\PY{p}{,}\PY{+w}{ }\PY{n}{β}\PY{p}{,}\PY{+w}{ }\PY{n}{U}\PY{p}{,}\PY{+w}{ }\PY{n}{V}\PY{p}{,}\PY{+w}{ }\PY{n}{X}\PY{p}{,}\PY{+w}{ }\PY{n}{E}\PY{+w}{ }\PY{o}{=}\PY{+w}{ }\PY{n}{ReconstrucaoLSQR}\PY{p}{(}\PY{n}{A}\PY{p}{,}\PY{n}{lena\PYZus{}vec\PYZus{}bn}\PY{p}{,}\PY{+w}{ }\PY{n}{δ}\PY{p}{,}\PY{+w}{ }\PY{l+m+mi}{10}\PY{p}{,}\PY{+w}{ }\PY{n}{τ}\PY{p}{,}\PY{+w}{ }\PY{n}{ortog}\PY{o}{=}\PY{n+nb}{true}\PY{p}{)}\PY{p}{;}
\end{Verbatim}
\end{tcolorbox}

    \begin{Verbatim}[commandchars=\\\{\}]
\textcolor{ansi-cyan-intense}{\textbf{[ }}\textcolor{ansi-cyan-intense}{\textbf{Info: }}Critério de discrepância aceito
com k=3
    \end{Verbatim}

    \begin{Verbatim}[commandchars=\\\{\}]
  2.643614 seconds (2.02 M allocations: 145.847 MiB, 28.64\% gc time, 70.70\%
compilation time)
    \end{Verbatim}

    Evolução da solução

    \begin{tcolorbox}[breakable, size=fbox, boxrule=1pt, pad at break*=1mm,colback=cellbackground, colframe=cellborder]
\prompt{In}{incolor}{ }{\boxspacing}
\begin{Verbatim}[commandchars=\\\{\}]
\PY{n}{mosaicview}\PY{p}{(}\PY{n}{vetor2imagem}\PY{o}{.}\PY{p}{(}\PY{n}{X}\PY{p}{)}\PY{o}{...}\PY{p}{;}\PY{+w}{ }\PY{n}{nrow}\PY{o}{=}\PY{l+m+mi}{1}\PY{p}{,}\PY{+w}{ }\PY{n}{npad}\PY{o}{=}\PY{l+m+mi}{3}\PY{p}{)}
\end{Verbatim}
\end{tcolorbox}
 
            
\prompt{Out}{outcolor}{ }{}
    
    \begin{center}
    \adjustimage{max size={0.9\linewidth}{0.9\paperheight}}{Trabalho_final_ALC2_files/Trabalho_final_ALC2_36_0.png}
    \end{center}
    { \hspace*{\fill} \\}
    

    Comparando lado a lado: Solução, imagem borrada e a imagem original

    \begin{tcolorbox}[breakable, size=fbox, boxrule=1pt, pad at break*=1mm,colback=cellbackground, colframe=cellborder]
\prompt{In}{incolor}{ }{\boxspacing}
\begin{Verbatim}[commandchars=\\\{\}]
\PY{n}{lena\PYZus{}rec\PYZus{}5\PYZus{}CO\PYZus{}11}\PY{+w}{ }\PY{o}{=}\PY{+w}{ }\PY{n}{vetor2imagem}\PY{p}{(}\PY{n}{X}\PY{p}{[}\PY{k}{end}\PY{p}{]}\PY{p}{)}

\PY{n}{mosaicview}\PY{p}{(}\PY{n}{lena\PYZus{}rec\PYZus{}5\PYZus{}CO\PYZus{}11}\PY{p}{,}\PY{+w}{ }\PY{n}{vetor2imagem}\PY{p}{(}\PY{n}{lena\PYZus{}vec\PYZus{}bn}\PY{p}{)}\PY{p}{,}\PY{+w}{ }\PY{n}{lena\PYZus{}img}\PY{p}{;}\PY{+w}{ }\PY{n}{nrow}\PY{o}{=}\PY{l+m+mi}{1}\PY{p}{,}\PY{+w}{ }\PY{n}{npad}\PY{o}{=}\PY{l+m+mi}{3}\PY{p}{)}
\end{Verbatim}
\end{tcolorbox}
 
            
\prompt{Out}{outcolor}{ }{}
    
    \begin{center}
    \adjustimage{max size={0.9\linewidth}{0.9\paperheight}}{Trabalho_final_ALC2_files/Trabalho_final_ALC2_38_0.png}
    \end{center}
    { \hspace*{\fill} \\}
    

    \hypertarget{erro-5-tau1.1-sem-ortogonalizauxe7uxe3o}{%
\subsubsection*{\texorpdfstring{Erro 5\%, \(\tau=1.1\), sem
ortogonalização}{Erro 5\%, \textbackslash tau=1.1, sem ortogonalização}}\label{erro-5-tau1.1-sem-ortogonalizauxe7uxe3o}}

    \begin{tcolorbox}[breakable, size=fbox, boxrule=1pt, pad at break*=1mm,colback=cellbackground, colframe=cellborder]
\prompt{In}{incolor}{ }{\boxspacing}
\begin{Verbatim}[commandchars=\\\{\}]
\PY{n}{α}\PY{p}{,}\PY{+w}{ }\PY{n}{β}\PY{p}{,}\PY{+w}{ }\PY{n}{U}\PY{p}{,}\PY{+w}{ }\PY{n}{V}\PY{p}{,}\PY{+w}{ }\PY{n}{X}\PY{p}{,}\PY{+w}{ }\PY{n}{E}\PY{+w}{ }\PY{o}{=}\PY{+w}{ }\PY{n}{ReconstrucaoLSQR}\PY{p}{(}\PY{n}{A}\PY{p}{,}\PY{n}{lena\PYZus{}vec\PYZus{}bn}\PY{p}{,}\PY{+w}{ }\PY{n}{δ}\PY{p}{,}\PY{+w}{ }\PY{l+m+mi}{10}\PY{p}{,}\PY{+w}{ }\PY{n}{τ}\PY{p}{,}\PY{+w}{ }\PY{n}{ortog}\PY{o}{=}\PY{n+nb}{false}\PY{p}{)}\PY{p}{;}
\end{Verbatim}
\end{tcolorbox}

    \begin{Verbatim}[commandchars=\\\{\}]
\textcolor{ansi-cyan-intense}{\textbf{[ }}\textcolor{ansi-cyan-intense}{\textbf{Info: }}Critério de discrepância aceito
com k=3
    \end{Verbatim}

    \begin{tcolorbox}[breakable, size=fbox, boxrule=1pt, pad at break*=1mm,colback=cellbackground, colframe=cellborder]
\prompt{In}{incolor}{ }{\boxspacing}
\begin{Verbatim}[commandchars=\\\{\}]
\PY{n}{mosaicview}\PY{p}{(}\PY{n}{vetor2imagem}\PY{o}{.}\PY{p}{(}\PY{n}{X}\PY{p}{)}\PY{o}{...}\PY{p}{;}\PY{+w}{ }\PY{n}{nrow}\PY{o}{=}\PY{l+m+mi}{1}\PY{p}{,}\PY{+w}{ }\PY{n}{npad}\PY{o}{=}\PY{l+m+mi}{3}\PY{p}{)}
\end{Verbatim}
\end{tcolorbox}
 
            
\prompt{Out}{outcolor}{ }{}
    
    \begin{center}
    \adjustimage{max size={0.9\linewidth}{0.9\paperheight}}{Trabalho_final_ALC2_files/Trabalho_final_ALC2_41_0.png}
    \end{center}
    { \hspace*{\fill} \\}
    

    Comparando lado a lado: Solução, imagem borrada e a imagem original

    \begin{tcolorbox}[breakable, size=fbox, boxrule=1pt, pad at break*=1mm,colback=cellbackground, colframe=cellborder]
\prompt{In}{incolor}{ }{\boxspacing}
\begin{Verbatim}[commandchars=\\\{\}]
\PY{n}{lena\PYZus{}rec\PYZus{}5\PYZus{}SO\PYZus{}11}\PY{+w}{ }\PY{o}{=}\PY{+w}{ }\PY{n}{vetor2imagem}\PY{p}{(}\PY{n}{X}\PY{p}{[}\PY{k}{end}\PY{p}{]}\PY{p}{)}
\PY{n}{mosaicview}\PY{p}{(}\PY{n}{lena\PYZus{}rec\PYZus{}5\PYZus{}SO\PYZus{}11}\PY{p}{,}\PY{+w}{ }\PY{n}{vetor2imagem}\PY{p}{(}\PY{n}{lena\PYZus{}vec\PYZus{}bn}\PY{p}{)}\PY{p}{,}\PY{+w}{ }\PY{n}{lena\PYZus{}img}\PY{p}{;}\PY{+w}{ }\PY{n}{nrow}\PY{o}{=}\PY{l+m+mi}{1}\PY{p}{,}\PY{+w}{ }\PY{n}{npad}\PY{o}{=}\PY{l+m+mi}{3}\PY{p}{)}
\end{Verbatim}
\end{tcolorbox}
 
            
\prompt{Out}{outcolor}{ }{}
    
    \begin{center}
    \adjustimage{max size={0.9\linewidth}{0.9\paperheight}}{Trabalho_final_ALC2_files/Trabalho_final_ALC2_43_0.png}
    \end{center}
    { \hspace*{\fill} \\}
    

    \hypertarget{erro-5-tau1.05-com-ortogonalizauxe7uxe3o}{%
\subsubsection*{\texorpdfstring{Erro 5\%, \(\tau=1.05\), com
ortogonalização}{Erro 5\%, \textbackslash tau=1.05, com ortogonalização}}\label{erro-5-tau1.05-com-ortogonalizauxe7uxe3o}}

    \begin{tcolorbox}[breakable, size=fbox, boxrule=1pt, pad at break*=1mm,colback=cellbackground, colframe=cellborder]
\prompt{In}{incolor}{ }{\boxspacing}
\begin{Verbatim}[commandchars=\\\{\}]
\PY{n}{A}\PY{+w}{ }\PY{o}{=}\PY{+w}{  }\PY{n}{blur}\PY{p}{(}\PY{n}{N}\PY{p}{,}\PY{l+m+mi}{5}\PY{p}{,}\PY{l+m+mf}{3.5}\PY{p}{)}\PY{p}{;}
\PY{n}{lena\PYZus{}vec\PYZus{}b}\PY{+w}{ }\PY{o}{=}\PY{+w}{ }\PY{n}{A}\PY{o}{*}\PY{n}{lena\PYZus{}vec}\PY{p}{;}

\PY{n}{erro}\PY{+w}{ }\PY{o}{=}\PY{+w}{ }\PY{n}{randn}\PY{p}{(}\PY{n}{N}\PY{o}{\PYZca{}}\PY{l+m+mi}{2}\PY{p}{)}
\PY{n}{erro}\PY{+w}{ }\PY{o}{=}\PY{+w}{ }\PY{n}{erro}\PY{o}{/}\PY{n}{norm}\PY{p}{(}\PY{n}{erro}\PY{p}{)}

\PY{n}{e}\PY{+w}{ }\PY{o}{=}\PY{+w}{ }\PY{l+m+mf}{0.05}
\PY{n}{τ}\PY{+w}{ }\PY{o}{=}\PY{+w}{ }\PY{l+m+mf}{1.05}

\PY{n}{lena\PYZus{}vec\PYZus{}bn}\PY{+w}{ }\PY{o}{=}\PY{+w}{ }\PY{n}{lena\PYZus{}vec\PYZus{}b}\PY{+w}{ }\PY{o}{+}\PY{+w}{ }\PY{n}{e}\PY{o}{*}\PY{n}{erro}\PY{o}{*}\PY{n}{norm}\PY{p}{(}\PY{n}{lena\PYZus{}vec\PYZus{}b}\PY{p}{)}
\PY{n}{lena\PYZus{}img\PYZus{}bn}\PY{+w}{ }\PY{o}{=}\PY{+w}{ }\PY{n}{vetor2imagem}\PY{p}{(}\PY{n}{lena\PYZus{}vec\PYZus{}bn}\PY{p}{)}
\end{Verbatim}
\end{tcolorbox}
 
            
\prompt{Out}{outcolor}{ }{}
    
    \begin{center}
    \adjustimage{max size={0.9\linewidth}{0.9\paperheight}}{Trabalho_final_ALC2_files/Trabalho_final_ALC2_45_0.png}
    \end{center}
    { \hspace*{\fill} \\}
    

    Critério de discrepância.

    \begin{tcolorbox}[breakable, size=fbox, boxrule=1pt, pad at break*=1mm,colback=cellbackground, colframe=cellborder]
\prompt{In}{incolor}{ }{\boxspacing}
\begin{Verbatim}[commandchars=\\\{\}]
\PY{n}{δ}\PY{+w}{ }\PY{o}{=}\PY{+w}{ }\PY{n}{norm}\PY{p}{(}\PY{n}{lena\PYZus{}vec\PYZus{}b}\PY{+w}{ }\PY{o}{\PYZhy{}}\PY{+w}{ }\PY{n}{lena\PYZus{}vec\PYZus{}bn}\PY{p}{)}
\end{Verbatim}
\end{tcolorbox}

            \begin{tcolorbox}[breakable, size=fbox, boxrule=.5pt, pad at break*=1mm, opacityfill=0]
\prompt{Out}{outcolor}{ }{\boxspacing}
\begin{Verbatim}[commandchars=\\\{\}]
4.21554656170367
\end{Verbatim}
\end{tcolorbox}
        
    \begin{tcolorbox}[breakable, size=fbox, boxrule=1pt, pad at break*=1mm,colback=cellbackground, colframe=cellborder]
\prompt{In}{incolor}{ }{\boxspacing}
\begin{Verbatim}[commandchars=\\\{\}]
\PY{n}{α}\PY{p}{,}\PY{+w}{ }\PY{n}{β}\PY{p}{,}\PY{+w}{ }\PY{n}{U}\PY{p}{,}\PY{+w}{ }\PY{n}{V}\PY{p}{,}\PY{+w}{ }\PY{n}{X}\PY{p}{,}\PY{+w}{ }\PY{n}{E}\PY{+w}{ }\PY{o}{=}\PY{+w}{ }\PY{n}{ReconstrucaoLSQR}\PY{p}{(}\PY{n}{A}\PY{p}{,}\PY{n}{lena\PYZus{}vec\PYZus{}bn}\PY{p}{,}\PY{+w}{ }\PY{n}{δ}\PY{p}{,}\PY{+w}{ }\PY{l+m+mi}{10}\PY{p}{,}\PY{+w}{ }\PY{n}{τ}\PY{p}{,}\PY{+w}{ }\PY{n}{ortog}\PY{o}{=}\PY{n+nb}{true}\PY{p}{)}\PY{p}{;}
\end{Verbatim}
\end{tcolorbox}

    \begin{Verbatim}[commandchars=\\\{\}]
\textcolor{ansi-cyan-intense}{\textbf{[ }}\textcolor{ansi-cyan-intense}{\textbf{Info: }}Critério de discrepância aceito
com k=4
    \end{Verbatim}

    Evolução da solução

    \begin{tcolorbox}[breakable, size=fbox, boxrule=1pt, pad at break*=1mm,colback=cellbackground, colframe=cellborder]
\prompt{In}{incolor}{ }{\boxspacing}
\begin{Verbatim}[commandchars=\\\{\}]
\PY{n}{mosaicview}\PY{p}{(}\PY{n}{vetor2imagem}\PY{o}{.}\PY{p}{(}\PY{n}{X}\PY{p}{)}\PY{o}{...}\PY{p}{;}\PY{+w}{ }\PY{n}{nrow}\PY{o}{=}\PY{l+m+mi}{1}\PY{p}{,}\PY{+w}{ }\PY{n}{npad}\PY{o}{=}\PY{l+m+mi}{3}\PY{p}{)}
\end{Verbatim}
\end{tcolorbox}
 
            
\prompt{Out}{outcolor}{ }{}
    
    \begin{center}
    \adjustimage{max size={0.9\linewidth}{0.9\paperheight}}{Trabalho_final_ALC2_files/Trabalho_final_ALC2_50_0.png}
    \end{center}
    { \hspace*{\fill} \\}
    

    Comparando lado a lado: Solução, imagem borrada e a imagem original

    \begin{tcolorbox}[breakable, size=fbox, boxrule=1pt, pad at break*=1mm,colback=cellbackground, colframe=cellborder]
\prompt{In}{incolor}{ }{\boxspacing}
\begin{Verbatim}[commandchars=\\\{\}]
\PY{n}{lena\PYZus{}rec\PYZus{}5\PYZus{}CO\PYZus{}105}\PY{+w}{ }\PY{o}{=}\PY{+w}{ }\PY{n}{vetor2imagem}\PY{p}{(}\PY{n}{X}\PY{p}{[}\PY{k}{end}\PY{p}{]}\PY{p}{)}

\PY{n}{mosaicview}\PY{p}{(}\PY{n}{lena\PYZus{}rec\PYZus{}5\PYZus{}CO\PYZus{}105}\PY{p}{,}\PY{+w}{ }\PY{n}{vetor2imagem}\PY{p}{(}\PY{n}{lena\PYZus{}vec\PYZus{}bn}\PY{p}{)}\PY{p}{,}\PY{+w}{ }\PY{n}{lena\PYZus{}img}\PY{p}{;}\PY{+w}{ }\PY{n}{nrow}\PY{o}{=}\PY{l+m+mi}{1}\PY{p}{,}\PY{+w}{ }\PY{n}{npad}\PY{o}{=}\PY{l+m+mi}{3}\PY{p}{)}
\end{Verbatim}
\end{tcolorbox}
 
            
\prompt{Out}{outcolor}{ }{}
    
    \begin{center}
    \adjustimage{max size={0.9\linewidth}{0.9\paperheight}}{Trabalho_final_ALC2_files/Trabalho_final_ALC2_52_0.png}
    \end{center}
    { \hspace*{\fill} \\}
    

    \hypertarget{erro-5-tau1.05-sem-ortogonalizauxe7uxe3o}{%
\subsubsection*{\texorpdfstring{Erro 5\%, \(\tau=1.05\), sem
ortogonalização}{Erro 5\%, \textbackslash tau=1.05, sem ortogonalização}}\label{erro-5-tau1.05-sem-ortogonalizauxe7uxe3o}}

    \begin{tcolorbox}[breakable, size=fbox, boxrule=1pt, pad at break*=1mm,colback=cellbackground, colframe=cellborder]
\prompt{In}{incolor}{ }{\boxspacing}
\begin{Verbatim}[commandchars=\\\{\}]
\PY{n}{α}\PY{p}{,}\PY{+w}{ }\PY{n}{β}\PY{p}{,}\PY{+w}{ }\PY{n}{U}\PY{p}{,}\PY{+w}{ }\PY{n}{V}\PY{p}{,}\PY{+w}{ }\PY{n}{X}\PY{p}{,}\PY{+w}{ }\PY{n}{E}\PY{+w}{ }\PY{o}{=}\PY{+w}{ }\PY{n}{ReconstrucaoLSQR}\PY{p}{(}\PY{n}{A}\PY{p}{,}\PY{n}{lena\PYZus{}vec\PYZus{}bn}\PY{p}{,}\PY{+w}{ }\PY{n}{δ}\PY{p}{,}\PY{+w}{ }\PY{l+m+mi}{10}\PY{p}{,}\PY{+w}{ }\PY{n}{τ}\PY{p}{,}\PY{+w}{ }\PY{n}{ortog}\PY{o}{=}\PY{n+nb}{false}\PY{p}{)}\PY{p}{;}
\end{Verbatim}
\end{tcolorbox}

    \begin{Verbatim}[commandchars=\\\{\}]
\textcolor{ansi-cyan-intense}{\textbf{[ }}\textcolor{ansi-cyan-intense}{\textbf{Info: }}Critério de discrepância aceito
com k=4
    \end{Verbatim}

    \begin{tcolorbox}[breakable, size=fbox, boxrule=1pt, pad at break*=1mm,colback=cellbackground, colframe=cellborder]
\prompt{In}{incolor}{ }{\boxspacing}
\begin{Verbatim}[commandchars=\\\{\}]
\PY{n}{mosaicview}\PY{p}{(}\PY{n}{vetor2imagem}\PY{o}{.}\PY{p}{(}\PY{n}{X}\PY{p}{)}\PY{o}{...}\PY{p}{;}\PY{+w}{ }\PY{n}{nrow}\PY{o}{=}\PY{l+m+mi}{1}\PY{p}{,}\PY{+w}{ }\PY{n}{npad}\PY{o}{=}\PY{l+m+mi}{3}\PY{p}{)}
\end{Verbatim}
\end{tcolorbox}
 
            
\prompt{Out}{outcolor}{ }{}
    
    \begin{center}
    \adjustimage{max size={0.9\linewidth}{0.9\paperheight}}{Trabalho_final_ALC2_files/Trabalho_final_ALC2_55_0.png}
    \end{center}
    { \hspace*{\fill} \\}
    

    Comparando lado a lado: Solução, imagem borrada e a imagem original

    \begin{tcolorbox}[breakable, size=fbox, boxrule=1pt, pad at break*=1mm,colback=cellbackground, colframe=cellborder]
\prompt{In}{incolor}{ }{\boxspacing}
\begin{Verbatim}[commandchars=\\\{\}]
\PY{n}{lena\PYZus{}rec\PYZus{}5\PYZus{}SO\PYZus{}105}\PY{+w}{ }\PY{o}{=}\PY{+w}{ }\PY{n}{vetor2imagem}\PY{p}{(}\PY{n}{X}\PY{p}{[}\PY{k}{end}\PY{p}{]}\PY{p}{)}
\PY{n}{mosaicview}\PY{p}{(}\PY{n}{lena\PYZus{}rec\PYZus{}5\PYZus{}SO\PYZus{}105}\PY{p}{,}\PY{+w}{ }\PY{n}{vetor2imagem}\PY{p}{(}\PY{n}{lena\PYZus{}vec\PYZus{}bn}\PY{p}{)}\PY{p}{,}\PY{+w}{ }\PY{n}{lena\PYZus{}img}\PY{p}{;}\PY{+w}{ }\PY{n}{nrow}\PY{o}{=}\PY{l+m+mi}{1}\PY{p}{,}\PY{+w}{ }\PY{n}{npad}\PY{o}{=}\PY{l+m+mi}{3}\PY{p}{)}
\end{Verbatim}
\end{tcolorbox}
 
            
\prompt{Out}{outcolor}{ }{}
    
    \begin{center}
    \adjustimage{max size={0.9\linewidth}{0.9\paperheight}}{Trabalho_final_ALC2_files/Trabalho_final_ALC2_57_0.png}
    \end{center}
    { \hspace*{\fill} \\}
    

    \hypertarget{erro-5-tau1.01-com-ortogonalizauxe7uxe3o}{%
\subsubsection*{\texorpdfstring{Erro 5\%, \(\tau=1.01\), com
ortogonalização}{Erro 5\%, \textbackslash tau=1.01, com ortogonalização}}\label{erro-5-tau1.01-com-ortogonalizauxe7uxe3o}}

    \begin{tcolorbox}[breakable, size=fbox, boxrule=1pt, pad at break*=1mm,colback=cellbackground, colframe=cellborder]
\prompt{In}{incolor}{ }{\boxspacing}
\begin{Verbatim}[commandchars=\\\{\}]
\PY{n}{A}\PY{+w}{ }\PY{o}{=}\PY{+w}{  }\PY{n}{blur}\PY{p}{(}\PY{n}{N}\PY{p}{,}\PY{l+m+mi}{5}\PY{p}{,}\PY{l+m+mf}{3.5}\PY{p}{)}\PY{p}{;}
\PY{n}{lena\PYZus{}vec\PYZus{}b}\PY{+w}{ }\PY{o}{=}\PY{+w}{ }\PY{n}{A}\PY{o}{*}\PY{n}{lena\PYZus{}vec}\PY{p}{;}

\PY{n}{erro}\PY{+w}{ }\PY{o}{=}\PY{+w}{ }\PY{n}{randn}\PY{p}{(}\PY{n}{N}\PY{o}{\PYZca{}}\PY{l+m+mi}{2}\PY{p}{)}
\PY{n}{erro}\PY{+w}{ }\PY{o}{=}\PY{+w}{ }\PY{n}{erro}\PY{o}{/}\PY{n}{norm}\PY{p}{(}\PY{n}{erro}\PY{p}{)}

\PY{n}{e}\PY{+w}{ }\PY{o}{=}\PY{+w}{ }\PY{l+m+mf}{0.05}
\PY{n}{τ}\PY{+w}{ }\PY{o}{=}\PY{+w}{ }\PY{l+m+mf}{1.01}

\PY{n}{lena\PYZus{}vec\PYZus{}bn}\PY{+w}{ }\PY{o}{=}\PY{+w}{ }\PY{n}{lena\PYZus{}vec\PYZus{}b}\PY{+w}{ }\PY{o}{+}\PY{+w}{ }\PY{n}{e}\PY{o}{*}\PY{n}{erro}\PY{o}{*}\PY{n}{norm}\PY{p}{(}\PY{n}{lena\PYZus{}vec\PYZus{}b}\PY{p}{)}
\PY{n}{lena\PYZus{}img\PYZus{}bn}\PY{+w}{ }\PY{o}{=}\PY{+w}{ }\PY{n}{vetor2imagem}\PY{p}{(}\PY{n}{lena\PYZus{}vec\PYZus{}bn}\PY{p}{)}
\end{Verbatim}
\end{tcolorbox}
 
            
\prompt{Out}{outcolor}{ }{}
    
    \begin{center}
    \adjustimage{max size={0.9\linewidth}{0.9\paperheight}}{Trabalho_final_ALC2_files/Trabalho_final_ALC2_59_0.png}
    \end{center}
    { \hspace*{\fill} \\}
    

    Critério de discrepância.

    \begin{tcolorbox}[breakable, size=fbox, boxrule=1pt, pad at break*=1mm,colback=cellbackground, colframe=cellborder]
\prompt{In}{incolor}{ }{\boxspacing}
\begin{Verbatim}[commandchars=\\\{\}]
\PY{n}{δ}\PY{+w}{ }\PY{o}{=}\PY{+w}{ }\PY{n}{norm}\PY{p}{(}\PY{n}{lena\PYZus{}vec\PYZus{}b}\PY{+w}{ }\PY{o}{\PYZhy{}}\PY{+w}{ }\PY{n}{lena\PYZus{}vec\PYZus{}bn}\PY{p}{)}
\end{Verbatim}
\end{tcolorbox}

            \begin{tcolorbox}[breakable, size=fbox, boxrule=.5pt, pad at break*=1mm, opacityfill=0]
\prompt{Out}{outcolor}{ }{\boxspacing}
\begin{Verbatim}[commandchars=\\\{\}]
4.21554656170367
\end{Verbatim}
\end{tcolorbox}
        
    \begin{tcolorbox}[breakable, size=fbox, boxrule=1pt, pad at break*=1mm,colback=cellbackground, colframe=cellborder]
\prompt{In}{incolor}{ }{\boxspacing}
\begin{Verbatim}[commandchars=\\\{\}]
\PY{n}{α}\PY{p}{,}\PY{+w}{ }\PY{n}{β}\PY{p}{,}\PY{+w}{ }\PY{n}{U}\PY{p}{,}\PY{+w}{ }\PY{n}{V}\PY{p}{,}\PY{+w}{ }\PY{n}{X}\PY{p}{,}\PY{+w}{ }\PY{n}{E}\PY{+w}{ }\PY{o}{=}\PY{+w}{ }\PY{n}{ReconstrucaoLSQR}\PY{p}{(}\PY{n}{A}\PY{p}{,}\PY{n}{lena\PYZus{}vec\PYZus{}bn}\PY{p}{,}\PY{+w}{ }\PY{n}{δ}\PY{p}{,}\PY{+w}{ }\PY{l+m+mi}{10}\PY{p}{,}\PY{+w}{ }\PY{n}{τ}\PY{p}{,}\PY{+w}{ }\PY{n}{ortog}\PY{o}{=}\PY{n+nb}{true}\PY{p}{)}\PY{p}{;}
\end{Verbatim}
\end{tcolorbox}

    \begin{Verbatim}[commandchars=\\\{\}]
\textcolor{ansi-cyan-intense}{\textbf{[ }}\textcolor{ansi-cyan-intense}{\textbf{Info: }}Critério de discrepância aceito
com k=5
    \end{Verbatim}

    Evolução da solução

    \begin{tcolorbox}[breakable, size=fbox, boxrule=1pt, pad at break*=1mm,colback=cellbackground, colframe=cellborder]
\prompt{In}{incolor}{ }{\boxspacing}
\begin{Verbatim}[commandchars=\\\{\}]
\PY{n}{mosaicview}\PY{p}{(}\PY{n}{vetor2imagem}\PY{o}{.}\PY{p}{(}\PY{n}{X}\PY{p}{)}\PY{o}{...}\PY{p}{;}\PY{+w}{ }\PY{n}{nrow}\PY{o}{=}\PY{l+m+mi}{1}\PY{p}{,}\PY{+w}{ }\PY{n}{npad}\PY{o}{=}\PY{l+m+mi}{3}\PY{p}{)}
\end{Verbatim}
\end{tcolorbox}
 
            
\prompt{Out}{outcolor}{ }{}
    
    \begin{center}
    \adjustimage{max size={0.9\linewidth}{0.9\paperheight}}{Trabalho_final_ALC2_files/Trabalho_final_ALC2_64_0.png}
    \end{center}
    { \hspace*{\fill} \\}
    

    Comparando lado a lado: Solução, imagem borrada e a imagem original

    \begin{tcolorbox}[breakable, size=fbox, boxrule=1pt, pad at break*=1mm,colback=cellbackground, colframe=cellborder]
\prompt{In}{incolor}{ }{\boxspacing}
\begin{Verbatim}[commandchars=\\\{\}]
\PY{n}{lena\PYZus{}rec\PYZus{}5\PYZus{}CO\PYZus{}101}\PY{+w}{ }\PY{o}{=}\PY{+w}{ }\PY{n}{vetor2imagem}\PY{p}{(}\PY{n}{X}\PY{p}{[}\PY{k}{end}\PY{p}{]}\PY{p}{)}

\PY{n}{mosaicview}\PY{p}{(}\PY{n}{lena\PYZus{}rec\PYZus{}5\PYZus{}CO\PYZus{}101}\PY{p}{,}\PY{+w}{ }\PY{n}{vetor2imagem}\PY{p}{(}\PY{n}{lena\PYZus{}vec\PYZus{}bn}\PY{p}{)}\PY{p}{,}\PY{+w}{ }\PY{n}{lena\PYZus{}img}\PY{p}{;}\PY{+w}{ }\PY{n}{nrow}\PY{o}{=}\PY{l+m+mi}{1}\PY{p}{,}\PY{+w}{ }\PY{n}{npad}\PY{o}{=}\PY{l+m+mi}{3}\PY{p}{)}
\end{Verbatim}
\end{tcolorbox}
 
            
\prompt{Out}{outcolor}{ }{}
    
    \begin{center}
    \adjustimage{max size={0.9\linewidth}{0.9\paperheight}}{Trabalho_final_ALC2_files/Trabalho_final_ALC2_66_0.png}
    \end{center}
    { \hspace*{\fill} \\}
    

    \hypertarget{erro-10-tau1.1-com-ortogonalizauxe7uxe3o}{%
\subsubsection*{\texorpdfstring{Erro 10\%, \(\tau=1.1\), com
ortogonalização}{Erro 10\%, \textbackslash tau=1.1, com ortogonalização}}\label{erro-10-tau1.1-com-ortogonalizauxe7uxe3o}}

    \begin{tcolorbox}[breakable, size=fbox, boxrule=1pt, pad at break*=1mm,colback=cellbackground, colframe=cellborder]
\prompt{In}{incolor}{ }{\boxspacing}
\begin{Verbatim}[commandchars=\\\{\}]
\PY{n}{A}\PY{+w}{ }\PY{o}{=}\PY{+w}{  }\PY{n}{blur}\PY{p}{(}\PY{n}{N}\PY{p}{,}\PY{l+m+mi}{5}\PY{p}{,}\PY{l+m+mf}{3.5}\PY{p}{)}\PY{p}{;}
\PY{n}{lena\PYZus{}vec\PYZus{}b}\PY{+w}{ }\PY{o}{=}\PY{+w}{ }\PY{n}{A}\PY{o}{*}\PY{n}{lena\PYZus{}vec}\PY{p}{;}

\PY{n}{erro}\PY{+w}{ }\PY{o}{=}\PY{+w}{ }\PY{n}{randn}\PY{p}{(}\PY{n}{N}\PY{o}{\PYZca{}}\PY{l+m+mi}{2}\PY{p}{)}
\PY{n}{erro}\PY{+w}{ }\PY{o}{=}\PY{+w}{ }\PY{n}{erro}\PY{o}{/}\PY{n}{norm}\PY{p}{(}\PY{n}{erro}\PY{p}{)}

\PY{n}{e}\PY{+w}{ }\PY{o}{=}\PY{+w}{ }\PY{l+m+mf}{0.10}
\PY{n}{τ}\PY{+w}{ }\PY{o}{=}\PY{+w}{ }\PY{l+m+mf}{1.1}

\PY{n}{lena\PYZus{}vec\PYZus{}bn}\PY{+w}{ }\PY{o}{=}\PY{+w}{ }\PY{n}{lena\PYZus{}vec\PYZus{}b}\PY{+w}{ }\PY{o}{+}\PY{+w}{ }\PY{n}{e}\PY{o}{*}\PY{n}{erro}\PY{o}{*}\PY{n}{norm}\PY{p}{(}\PY{n}{lena\PYZus{}vec\PYZus{}b}\PY{p}{)}
\PY{n}{lena\PYZus{}img\PYZus{}bn}\PY{+w}{ }\PY{o}{=}\PY{+w}{ }\PY{n}{vetor2imagem}\PY{p}{(}\PY{n}{lena\PYZus{}vec\PYZus{}bn}\PY{p}{)}
\end{Verbatim}
\end{tcolorbox}
 
            
\prompt{Out}{outcolor}{ }{}
    
    \begin{center}
    \adjustimage{max size={0.9\linewidth}{0.9\paperheight}}{Trabalho_final_ALC2_files/Trabalho_final_ALC2_68_0.png}
    \end{center}
    { \hspace*{\fill} \\}
    

    Critério de discrepância.

    \begin{tcolorbox}[breakable, size=fbox, boxrule=1pt, pad at break*=1mm,colback=cellbackground, colframe=cellborder]
\prompt{In}{incolor}{ }{\boxspacing}
\begin{Verbatim}[commandchars=\\\{\}]
\PY{n}{δ}\PY{+w}{ }\PY{o}{=}\PY{+w}{ }\PY{n}{norm}\PY{p}{(}\PY{n}{lena\PYZus{}vec\PYZus{}b}\PY{+w}{ }\PY{o}{\PYZhy{}}\PY{+w}{ }\PY{n}{lena\PYZus{}vec\PYZus{}bn}\PY{p}{)}
\end{Verbatim}
\end{tcolorbox}

            \begin{tcolorbox}[breakable, size=fbox, boxrule=.5pt, pad at break*=1mm, opacityfill=0]
\prompt{Out}{outcolor}{ }{\boxspacing}
\begin{Verbatim}[commandchars=\\\{\}]
8.43109312340734
\end{Verbatim}
\end{tcolorbox}
        
    \begin{tcolorbox}[breakable, size=fbox, boxrule=1pt, pad at break*=1mm,colback=cellbackground, colframe=cellborder]
\prompt{In}{incolor}{ }{\boxspacing}
\begin{Verbatim}[commandchars=\\\{\}]
\PY{n+nd}{@time}\PY{+w}{ }\PY{n}{α}\PY{p}{,}\PY{+w}{ }\PY{n}{β}\PY{p}{,}\PY{+w}{ }\PY{n}{U}\PY{p}{,}\PY{+w}{ }\PY{n}{V}\PY{p}{,}\PY{+w}{ }\PY{n}{X}\PY{p}{,}\PY{+w}{ }\PY{n}{E}\PY{+w}{ }\PY{o}{=}\PY{+w}{ }\PY{n}{ReconstrucaoLSQR}\PY{p}{(}\PY{n}{A}\PY{p}{,}\PY{n}{lena\PYZus{}vec\PYZus{}bn}\PY{p}{,}\PY{+w}{ }\PY{n}{δ}\PY{p}{,}\PY{+w}{ }\PY{l+m+mi}{10}\PY{p}{,}\PY{+w}{ }\PY{n}{τ}\PY{p}{,}\PY{+w}{ }\PY{n}{ortog}\PY{o}{=}\PY{n+nb}{true}\PY{p}{)}\PY{p}{;}
\end{Verbatim}
\end{tcolorbox}

    \begin{Verbatim}[commandchars=\\\{\}]
\textcolor{ansi-cyan-intense}{\textbf{[ }}\textcolor{ansi-cyan-intense}{\textbf{Info: }}Critério de discrepância aceito
com k=2
    \end{Verbatim}

    \begin{Verbatim}[commandchars=\\\{\}]
  0.079342 seconds (411 allocations: 27.081 MiB, 4.25\% gc time)
    \end{Verbatim}

    Evolução da solução

    \begin{tcolorbox}[breakable, size=fbox, boxrule=1pt, pad at break*=1mm,colback=cellbackground, colframe=cellborder]
\prompt{In}{incolor}{ }{\boxspacing}
\begin{Verbatim}[commandchars=\\\{\}]
\PY{n}{mosaicview}\PY{p}{(}\PY{n}{vetor2imagem}\PY{o}{.}\PY{p}{(}\PY{n}{X}\PY{p}{)}\PY{o}{...}\PY{p}{;}\PY{+w}{ }\PY{n}{nrow}\PY{o}{=}\PY{l+m+mi}{1}\PY{p}{,}\PY{+w}{ }\PY{n}{npad}\PY{o}{=}\PY{l+m+mi}{3}\PY{p}{)}
\end{Verbatim}
\end{tcolorbox}
 
            
\prompt{Out}{outcolor}{ }{}
    
    \begin{center}
    \adjustimage{max size={0.9\linewidth}{0.9\paperheight}}{Trabalho_final_ALC2_files/Trabalho_final_ALC2_73_0.png}
    \end{center}
    { \hspace*{\fill} \\}
    

    Comparando lado a lado: Solução, imagem borrada e a imagem original

    \begin{tcolorbox}[breakable, size=fbox, boxrule=1pt, pad at break*=1mm,colback=cellbackground, colframe=cellborder]
\prompt{In}{incolor}{ }{\boxspacing}
\begin{Verbatim}[commandchars=\\\{\}]
\PY{n}{lena\PYZus{}rec\PYZus{}10\PYZus{}CO\PYZus{}11}\PY{+w}{ }\PY{o}{=}\PY{+w}{ }\PY{n}{vetor2imagem}\PY{p}{(}\PY{n}{X}\PY{p}{[}\PY{k}{end}\PY{p}{]}\PY{p}{)}

\PY{n}{mosaicview}\PY{p}{(}\PY{n}{lena\PYZus{}rec\PYZus{}10\PYZus{}CO\PYZus{}11}\PY{p}{,}\PY{+w}{ }\PY{n}{vetor2imagem}\PY{p}{(}\PY{n}{lena\PYZus{}vec\PYZus{}bn}\PY{p}{)}\PY{p}{,}\PY{+w}{ }\PY{n}{lena\PYZus{}img}\PY{p}{;}\PY{+w}{ }\PY{n}{nrow}\PY{o}{=}\PY{l+m+mi}{1}\PY{p}{,}\PY{+w}{ }\PY{n}{npad}\PY{o}{=}\PY{l+m+mi}{3}\PY{p}{)}
\end{Verbatim}
\end{tcolorbox}
 
            
\prompt{Out}{outcolor}{ }{}
    
    \begin{center}
    \adjustimage{max size={0.9\linewidth}{0.9\paperheight}}{Trabalho_final_ALC2_files/Trabalho_final_ALC2_75_0.png}
    \end{center}
    { \hspace*{\fill} \\}
    

    \hypertarget{erro-10-tau1.1-sem-ortogonalizauxe7uxe3o}{%
\subsubsection*{\texorpdfstring{Erro 10\%, \(\tau=1.1\), sem
ortogonalização}{Erro 10\%, \textbackslash tau=1.1, sem ortogonalização}}\label{erro-10-tau1.1-sem-ortogonalizauxe7uxe3o}}

    \begin{tcolorbox}[breakable, size=fbox, boxrule=1pt, pad at break*=1mm,colback=cellbackground, colframe=cellborder]
\prompt{In}{incolor}{ }{\boxspacing}
\begin{Verbatim}[commandchars=\\\{\}]
\PY{n+nd}{@time}\PY{+w}{ }\PY{n}{α}\PY{p}{,}\PY{+w}{ }\PY{n}{β}\PY{p}{,}\PY{+w}{ }\PY{n}{U}\PY{p}{,}\PY{+w}{ }\PY{n}{V}\PY{p}{,}\PY{+w}{ }\PY{n}{X}\PY{p}{,}\PY{+w}{ }\PY{n}{E}\PY{+w}{ }\PY{o}{=}\PY{+w}{ }\PY{n}{ReconstrucaoLSQR}\PY{p}{(}\PY{n}{A}\PY{p}{,}\PY{n}{lena\PYZus{}vec\PYZus{}bn}\PY{p}{,}\PY{+w}{ }\PY{n}{δ}\PY{p}{,}\PY{+w}{ }\PY{l+m+mi}{10}\PY{p}{,}\PY{+w}{ }\PY{n}{τ}\PY{p}{,}\PY{+w}{ }\PY{n}{ortog}\PY{o}{=}\PY{n+nb}{false}\PY{p}{)}\PY{p}{;}
\end{Verbatim}
\end{tcolorbox}

    \begin{Verbatim}[commandchars=\\\{\}]
\textcolor{ansi-cyan-intense}{\textbf{[ }}\textcolor{ansi-cyan-intense}{\textbf{Info: }}Critério de discrepância aceito
com k=2
    \end{Verbatim}

    \begin{Verbatim}[commandchars=\\\{\}]
  0.073843 seconds (351 allocations: 17.080 MiB)
    \end{Verbatim}

    \begin{tcolorbox}[breakable, size=fbox, boxrule=1pt, pad at break*=1mm,colback=cellbackground, colframe=cellborder]
\prompt{In}{incolor}{ }{\boxspacing}
\begin{Verbatim}[commandchars=\\\{\}]
\PY{n}{mosaicview}\PY{p}{(}\PY{n}{vetor2imagem}\PY{o}{.}\PY{p}{(}\PY{n}{X}\PY{p}{)}\PY{o}{...}\PY{p}{;}\PY{+w}{ }\PY{n}{nrow}\PY{o}{=}\PY{l+m+mi}{1}\PY{p}{,}\PY{+w}{ }\PY{n}{npad}\PY{o}{=}\PY{l+m+mi}{3}\PY{p}{)}
\end{Verbatim}
\end{tcolorbox}
 
            
\prompt{Out}{outcolor}{ }{}
    
    \begin{center}
    \adjustimage{max size={0.9\linewidth}{0.9\paperheight}}{Trabalho_final_ALC2_files/Trabalho_final_ALC2_78_0.png}
    \end{center}
    { \hspace*{\fill} \\}
    

    Comparando lado a lado: Solução, imagem borrada e a imagem original

    \begin{tcolorbox}[breakable, size=fbox, boxrule=1pt, pad at break*=1mm,colback=cellbackground, colframe=cellborder]
\prompt{In}{incolor}{ }{\boxspacing}
\begin{Verbatim}[commandchars=\\\{\}]
\PY{n}{lena\PYZus{}rec\PYZus{}10\PYZus{}SO\PYZus{}11}\PY{+w}{ }\PY{o}{=}\PY{+w}{ }\PY{n}{vetor2imagem}\PY{p}{(}\PY{n}{X}\PY{p}{[}\PY{k}{end}\PY{p}{]}\PY{p}{)}
\PY{n}{mosaicview}\PY{p}{(}\PY{n}{lena\PYZus{}rec\PYZus{}10\PYZus{}CO\PYZus{}11}\PY{p}{,}\PY{+w}{ }\PY{n}{vetor2imagem}\PY{p}{(}\PY{n}{lena\PYZus{}vec\PYZus{}bn}\PY{p}{)}\PY{p}{,}\PY{+w}{ }\PY{n}{lena\PYZus{}img}\PY{p}{;}\PY{+w}{ }\PY{n}{nrow}\PY{o}{=}\PY{l+m+mi}{1}\PY{p}{,}\PY{+w}{ }\PY{n}{npad}\PY{o}{=}\PY{l+m+mi}{3}\PY{p}{)}
\end{Verbatim}
\end{tcolorbox}
 
            
\prompt{Out}{outcolor}{ }{}
    
    \begin{center}
    \adjustimage{max size={0.9\linewidth}{0.9\paperheight}}{Trabalho_final_ALC2_files/Trabalho_final_ALC2_80_0.png}
    \end{center}
    { \hspace*{\fill} \\}
    

    \hypertarget{erro-15-tau1.05-com-ortogonalizauxe7uxe3o}{%
\subsubsection*{\texorpdfstring{Erro 15\%, \(\tau=1.05\), com
ortogonalização}{Erro 15\%, \textbackslash tau=1.05, com ortogonalização}}\label{erro-15-tau1.05-com-ortogonalizauxe7uxe3o}}

    \begin{tcolorbox}[breakable, size=fbox, boxrule=1pt, pad at break*=1mm,colback=cellbackground, colframe=cellborder]
\prompt{In}{incolor}{ }{\boxspacing}
\begin{Verbatim}[commandchars=\\\{\}]
\PY{n}{A}\PY{+w}{ }\PY{o}{=}\PY{+w}{  }\PY{n}{blur}\PY{p}{(}\PY{n}{N}\PY{p}{,}\PY{l+m+mi}{5}\PY{p}{,}\PY{l+m+mf}{3.5}\PY{p}{)}\PY{p}{;}
\PY{n}{lena\PYZus{}vec\PYZus{}b}\PY{+w}{ }\PY{o}{=}\PY{+w}{ }\PY{n}{A}\PY{o}{*}\PY{n}{lena\PYZus{}vec}\PY{p}{;}

\PY{n}{erro}\PY{+w}{ }\PY{o}{=}\PY{+w}{ }\PY{n}{randn}\PY{p}{(}\PY{n}{N}\PY{o}{\PYZca{}}\PY{l+m+mi}{2}\PY{p}{)}
\PY{n}{erro}\PY{+w}{ }\PY{o}{=}\PY{+w}{ }\PY{n}{erro}\PY{o}{/}\PY{n}{norm}\PY{p}{(}\PY{n}{erro}\PY{p}{)}

\PY{n}{e}\PY{+w}{ }\PY{o}{=}\PY{+w}{ }\PY{l+m+mf}{0.15}
\PY{n}{τ}\PY{+w}{ }\PY{o}{=}\PY{+w}{ }\PY{l+m+mf}{1.05}

\PY{n}{lena\PYZus{}vec\PYZus{}bn}\PY{+w}{ }\PY{o}{=}\PY{+w}{ }\PY{n}{lena\PYZus{}vec\PYZus{}b}\PY{+w}{ }\PY{o}{+}\PY{+w}{ }\PY{n}{e}\PY{o}{*}\PY{n}{erro}\PY{o}{*}\PY{n}{norm}\PY{p}{(}\PY{n}{lena\PYZus{}vec\PYZus{}b}\PY{p}{)}
\PY{n}{lena\PYZus{}img\PYZus{}bn}\PY{+w}{ }\PY{o}{=}\PY{+w}{ }\PY{n}{vetor2imagem}\PY{p}{(}\PY{n}{lena\PYZus{}vec\PYZus{}bn}\PY{p}{)}
\end{Verbatim}
\end{tcolorbox}
 
            
\prompt{Out}{outcolor}{ }{}
    
    \begin{center}
    \adjustimage{max size={0.9\linewidth}{0.9\paperheight}}{Trabalho_final_ALC2_files/Trabalho_final_ALC2_82_0.png}
    \end{center}
    { \hspace*{\fill} \\}
    

    Critério de discrepância.

    \begin{tcolorbox}[breakable, size=fbox, boxrule=1pt, pad at break*=1mm,colback=cellbackground, colframe=cellborder]
\prompt{In}{incolor}{ }{\boxspacing}
\begin{Verbatim}[commandchars=\\\{\}]
\PY{n}{δ}\PY{+w}{ }\PY{o}{=}\PY{+w}{ }\PY{n}{norm}\PY{p}{(}\PY{n}{lena\PYZus{}vec\PYZus{}b}\PY{+w}{ }\PY{o}{\PYZhy{}}\PY{+w}{ }\PY{n}{lena\PYZus{}vec\PYZus{}bn}\PY{p}{)}
\end{Verbatim}
\end{tcolorbox}

            \begin{tcolorbox}[breakable, size=fbox, boxrule=.5pt, pad at break*=1mm, opacityfill=0]
\prompt{Out}{outcolor}{ }{\boxspacing}
\begin{Verbatim}[commandchars=\\\{\}]
12.646639685111008
\end{Verbatim}
\end{tcolorbox}
        
    \begin{tcolorbox}[breakable, size=fbox, boxrule=1pt, pad at break*=1mm,colback=cellbackground, colframe=cellborder]
\prompt{In}{incolor}{ }{\boxspacing}
\begin{Verbatim}[commandchars=\\\{\}]
\PY{n+nd}{@time}\PY{+w}{ }\PY{n}{α}\PY{p}{,}\PY{+w}{ }\PY{n}{β}\PY{p}{,}\PY{+w}{ }\PY{n}{U}\PY{p}{,}\PY{+w}{ }\PY{n}{V}\PY{p}{,}\PY{+w}{ }\PY{n}{X}\PY{p}{,}\PY{+w}{ }\PY{n}{E}\PY{+w}{ }\PY{o}{=}\PY{+w}{ }\PY{n}{ReconstrucaoLSQR}\PY{p}{(}\PY{n}{A}\PY{p}{,}\PY{n}{lena\PYZus{}vec\PYZus{}bn}\PY{p}{,}\PY{+w}{ }\PY{n}{δ}\PY{p}{,}\PY{+w}{ }\PY{l+m+mi}{10}\PY{p}{,}\PY{+w}{ }\PY{n}{τ}\PY{p}{,}\PY{+w}{ }\PY{n}{ortog}\PY{o}{=}\PY{n+nb}{true}\PY{p}{)}\PY{p}{;}
\end{Verbatim}
\end{tcolorbox}

    \begin{Verbatim}[commandchars=\\\{\}]
\textcolor{ansi-cyan-intense}{\textbf{[ }}\textcolor{ansi-cyan-intense}{\textbf{Info: }}Critério de discrepância aceito
com k=2
    \end{Verbatim}

    \begin{Verbatim}[commandchars=\\\{\}]
  0.078629 seconds (411 allocations: 27.081 MiB)
    \end{Verbatim}

    Evolução da solução

    \begin{tcolorbox}[breakable, size=fbox, boxrule=1pt, pad at break*=1mm,colback=cellbackground, colframe=cellborder]
\prompt{In}{incolor}{ }{\boxspacing}
\begin{Verbatim}[commandchars=\\\{\}]
\PY{n}{mosaicview}\PY{p}{(}\PY{n}{vetor2imagem}\PY{o}{.}\PY{p}{(}\PY{n}{X}\PY{p}{)}\PY{o}{...}\PY{p}{;}\PY{+w}{ }\PY{n}{nrow}\PY{o}{=}\PY{l+m+mi}{1}\PY{p}{,}\PY{+w}{ }\PY{n}{npad}\PY{o}{=}\PY{l+m+mi}{3}\PY{p}{)}
\end{Verbatim}
\end{tcolorbox}
 
            
\prompt{Out}{outcolor}{ }{}
    
    \begin{center}
    \adjustimage{max size={0.9\linewidth}{0.9\paperheight}}{Trabalho_final_ALC2_files/Trabalho_final_ALC2_87_0.png}
    \end{center}
    { \hspace*{\fill} \\}
    

    Comparando lado a lado: Solução, imagem borrada e a imagem original

    \begin{tcolorbox}[breakable, size=fbox, boxrule=1pt, pad at break*=1mm,colback=cellbackground, colframe=cellborder]
\prompt{In}{incolor}{ }{\boxspacing}
\begin{Verbatim}[commandchars=\\\{\}]
\PY{n}{lena\PYZus{}rec\PYZus{}15\PYZus{}CO\PYZus{}105}\PY{+w}{ }\PY{o}{=}\PY{+w}{ }\PY{n}{vetor2imagem}\PY{p}{(}\PY{n}{X}\PY{p}{[}\PY{k}{end}\PY{p}{]}\PY{p}{)}

\PY{n}{mosaicview}\PY{p}{(}\PY{n}{lena\PYZus{}rec\PYZus{}15\PYZus{}CO\PYZus{}105}\PY{p}{,}\PY{+w}{ }\PY{n}{vetor2imagem}\PY{p}{(}\PY{n}{lena\PYZus{}vec\PYZus{}bn}\PY{p}{)}\PY{p}{,}\PY{+w}{ }\PY{n}{lena\PYZus{}img}\PY{p}{;}\PY{+w}{ }\PY{n}{nrow}\PY{o}{=}\PY{l+m+mi}{1}\PY{p}{,}\PY{+w}{ }\PY{n}{npad}\PY{o}{=}\PY{l+m+mi}{3}\PY{p}{)}
\end{Verbatim}
\end{tcolorbox}
 
            
\prompt{Out}{outcolor}{ }{}
    
    \begin{center}
    \adjustimage{max size={0.9\linewidth}{0.9\paperheight}}{Trabalho_final_ALC2_files/Trabalho_final_ALC2_89_0.png}
    \end{center}
    { \hspace*{\fill} \\}
    

    \begin{tcolorbox}[breakable, size=fbox, boxrule=1pt, pad at break*=1mm,colback=cellbackground, colframe=cellborder]
\prompt{In}{incolor}{ }{\boxspacing}
\begin{Verbatim}[commandchars=\\\{\}]
\PY{n}{A}\PY{+w}{ }\PY{o}{=}\PY{+w}{  }\PY{n}{blur}\PY{p}{(}\PY{n}{N}\PY{p}{,}\PY{l+m+mi}{5}\PY{p}{,}\PY{l+m+mf}{3.5}\PY{p}{)}\PY{p}{;}
\PY{n}{lena\PYZus{}vec\PYZus{}b}\PY{+w}{ }\PY{o}{=}\PY{+w}{ }\PY{n}{A}\PY{o}{*}\PY{n}{lena\PYZus{}vec}\PY{p}{;}

\PY{n}{erro}\PY{+w}{ }\PY{o}{=}\PY{+w}{ }\PY{n}{randn}\PY{p}{(}\PY{n}{N}\PY{o}{\PYZca{}}\PY{l+m+mi}{2}\PY{p}{)}
\PY{n}{erro}\PY{+w}{ }\PY{o}{=}\PY{+w}{ }\PY{n}{erro}\PY{o}{/}\PY{n}{norm}\PY{p}{(}\PY{n}{erro}\PY{p}{)}

\PY{n}{e}\PY{+w}{ }\PY{o}{=}\PY{+w}{ }\PY{l+m+mf}{0.05}
\PY{n}{τ}\PY{+w}{ }\PY{o}{=}\PY{+w}{ }\PY{l+m+mf}{1.05}

\PY{n}{lena\PYZus{}vec\PYZus{}bn}\PY{+w}{ }\PY{o}{=}\PY{+w}{ }\PY{n}{lena\PYZus{}vec\PYZus{}b}\PY{+w}{ }\PY{o}{+}\PY{+w}{ }\PY{n}{e}\PY{o}{*}\PY{n}{erro}\PY{o}{*}\PY{n}{norm}\PY{p}{(}\PY{n}{lena\PYZus{}vec\PYZus{}b}\PY{p}{)}
\PY{n}{lena\PYZus{}img\PYZus{}bn}\PY{+w}{ }\PY{o}{=}\PY{+w}{ }\PY{n}{vetor2imagem}\PY{p}{(}\PY{n}{lena\PYZus{}vec\PYZus{}bn}\PY{p}{)}
\end{Verbatim}
\end{tcolorbox}
 
            
\prompt{Out}{outcolor}{ }{}
    
    \begin{center}
    \adjustimage{max size={0.9\linewidth}{0.9\paperheight}}{Trabalho_final_ALC2_files/Trabalho_final_ALC2_90_0.png}
    \end{center}
    { \hspace*{\fill} \\}
    

    Critério de discrepância.

    \begin{tcolorbox}[breakable, size=fbox, boxrule=1pt, pad at break*=1mm,colback=cellbackground, colframe=cellborder]
\prompt{In}{incolor}{ }{\boxspacing}
\begin{Verbatim}[commandchars=\\\{\}]
\PY{n}{δ}\PY{+w}{ }\PY{o}{=}\PY{+w}{ }\PY{n}{norm}\PY{p}{(}\PY{n}{lena\PYZus{}vec\PYZus{}b}\PY{+w}{ }\PY{o}{\PYZhy{}}\PY{+w}{ }\PY{n}{lena\PYZus{}vec\PYZus{}bn}\PY{p}{)}
\end{Verbatim}
\end{tcolorbox}

            \begin{tcolorbox}[breakable, size=fbox, boxrule=.5pt, pad at break*=1mm, opacityfill=0]
\prompt{Out}{outcolor}{ }{\boxspacing}
\begin{Verbatim}[commandchars=\\\{\}]
4.21554656170367
\end{Verbatim}
\end{tcolorbox}
        
    \begin{tcolorbox}[breakable, size=fbox, boxrule=1pt, pad at break*=1mm,colback=cellbackground, colframe=cellborder]
\prompt{In}{incolor}{ }{\boxspacing}
\begin{Verbatim}[commandchars=\\\{\}]
\PY{n+nd}{@time}\PY{+w}{  }\PY{n}{α}\PY{p}{,}\PY{+w}{ }\PY{n}{β}\PY{p}{,}\PY{+w}{ }\PY{n}{U}\PY{p}{,}\PY{+w}{ }\PY{n}{V}\PY{p}{,}\PY{+w}{ }\PY{n}{X}\PY{p}{,}\PY{+w}{ }\PY{n}{E}\PY{+w}{ }\PY{o}{=}\PY{+w}{ }\PY{n}{ReconstrucaoLSQR}\PY{p}{(}\PY{n}{A}\PY{p}{,}\PY{n}{lena\PYZus{}vec\PYZus{}bn}\PY{p}{,}\PY{+w}{ }\PY{n}{δ}\PY{p}{,}\PY{+w}{ }\PY{l+m+mi}{10}\PY{p}{,}\PY{+w}{ }\PY{n}{τ}\PY{p}{,}\PY{+w}{ }\PY{n}{ortog}\PY{o}{=}\PY{n+nb}{true}\PY{p}{)}\PY{p}{;}
\end{Verbatim}
\end{tcolorbox}

    \begin{Verbatim}[commandchars=\\\{\}]
\textcolor{ansi-cyan-intense}{\textbf{[ }}\textcolor{ansi-cyan-intense}{\textbf{Info: }}Critério de discrepância aceito
com k=4
    \end{Verbatim}

    \begin{Verbatim}[commandchars=\\\{\}]
  0.176432 seconds (749 allocations: 61.661 MiB, 6.16\% gc time)
    \end{Verbatim}

    Evolução da solução

    \begin{tcolorbox}[breakable, size=fbox, boxrule=1pt, pad at break*=1mm,colback=cellbackground, colframe=cellborder]
\prompt{In}{incolor}{ }{\boxspacing}
\begin{Verbatim}[commandchars=\\\{\}]
\PY{n}{mosaicview}\PY{p}{(}\PY{n}{vetor2imagem}\PY{o}{.}\PY{p}{(}\PY{n}{X}\PY{p}{)}\PY{o}{...}\PY{p}{;}\PY{+w}{ }\PY{n}{nrow}\PY{o}{=}\PY{l+m+mi}{1}\PY{p}{,}\PY{+w}{ }\PY{n}{npad}\PY{o}{=}\PY{l+m+mi}{3}\PY{p}{)}
\end{Verbatim}
\end{tcolorbox}
 
            
\prompt{Out}{outcolor}{ }{}
    
    \begin{center}
    \adjustimage{max size={0.9\linewidth}{0.9\paperheight}}{Trabalho_final_ALC2_files/Trabalho_final_ALC2_95_0.png}
    \end{center}
    { \hspace*{\fill} \\}
    

    Comparando lado a lado: Solução, imagem borrada e a imagem original

    \begin{tcolorbox}[breakable, size=fbox, boxrule=1pt, pad at break*=1mm,colback=cellbackground, colframe=cellborder]
\prompt{In}{incolor}{ }{\boxspacing}
\begin{Verbatim}[commandchars=\\\{\}]
\PY{n}{lena\PYZus{}rec\PYZus{}5\PYZus{}CO\PYZus{}105}\PY{+w}{ }\PY{o}{=}\PY{+w}{ }\PY{n}{vetor2imagem}\PY{p}{(}\PY{n}{X}\PY{p}{[}\PY{k}{end}\PY{p}{]}\PY{p}{)}

\PY{n}{mosaicview}\PY{p}{(}\PY{n}{lena\PYZus{}rec\PYZus{}5\PYZus{}CO\PYZus{}105}\PY{p}{,}\PY{+w}{ }\PY{n}{vetor2imagem}\PY{p}{(}\PY{n}{lena\PYZus{}vec\PYZus{}bn}\PY{p}{)}\PY{p}{,}\PY{+w}{ }\PY{n}{lena\PYZus{}img}\PY{p}{;}\PY{+w}{ }\PY{n}{nrow}\PY{o}{=}\PY{l+m+mi}{1}\PY{p}{,}\PY{+w}{ }\PY{n}{npad}\PY{o}{=}\PY{l+m+mi}{3}\PY{p}{)}
\end{Verbatim}
\end{tcolorbox}
 
            
\prompt{Out}{outcolor}{ }{}
    
    \begin{center}
    \adjustimage{max size={0.9\linewidth}{0.9\paperheight}}{Trabalho_final_ALC2_files/Trabalho_final_ALC2_97_0.png}
    \end{center}
    { \hspace*{\fill} \\}
    

    \hypertarget{erro-10-tau0-com-ortogonalizauxe7uxe3o}{%
\subsubsection*{\texorpdfstring{Erro 10\%, \(\tau=0\), com
ortogonalização}{Erro 10\%, \textbackslash tau=0, com ortogonalização}}\label{erro-10-tau0-com-ortogonalizauxe7uxe3o}}

    Vamos calcular um exemplo sem o critério de parada, colocando
\(\tau=0\).

Nesse caso, após algumas iterações, o método começa a acumular muito
erro.

    \begin{tcolorbox}[breakable, size=fbox, boxrule=1pt, pad at break*=1mm,colback=cellbackground, colframe=cellborder]
\prompt{In}{incolor}{ }{\boxspacing}
\begin{Verbatim}[commandchars=\\\{\}]
\PY{n}{A}\PY{+w}{ }\PY{o}{=}\PY{+w}{  }\PY{n}{blur}\PY{p}{(}\PY{n}{N}\PY{p}{,}\PY{l+m+mi}{5}\PY{p}{,}\PY{l+m+mf}{3.5}\PY{p}{)}\PY{p}{;}
\PY{n}{lena\PYZus{}vec\PYZus{}b}\PY{+w}{ }\PY{o}{=}\PY{+w}{ }\PY{n}{A}\PY{o}{*}\PY{n}{lena\PYZus{}vec}\PY{p}{;}

\PY{n}{erro}\PY{+w}{ }\PY{o}{=}\PY{+w}{ }\PY{n}{randn}\PY{p}{(}\PY{n}{N}\PY{o}{\PYZca{}}\PY{l+m+mi}{2}\PY{p}{)}
\PY{n}{erro}\PY{+w}{ }\PY{o}{=}\PY{+w}{ }\PY{n}{erro}\PY{o}{/}\PY{n}{norm}\PY{p}{(}\PY{n}{erro}\PY{p}{)}

\PY{n}{e}\PY{+w}{ }\PY{o}{=}\PY{+w}{ }\PY{l+m+mf}{0.10}
\PY{n}{τ}\PY{+w}{ }\PY{o}{=}\PY{+w}{ }\PY{l+m+mi}{0}

\PY{n}{lena\PYZus{}vec\PYZus{}bn}\PY{+w}{ }\PY{o}{=}\PY{+w}{ }\PY{n}{lena\PYZus{}vec\PYZus{}b}\PY{+w}{ }\PY{o}{+}\PY{+w}{ }\PY{n}{e}\PY{o}{*}\PY{n}{erro}\PY{o}{*}\PY{n}{norm}\PY{p}{(}\PY{n}{lena\PYZus{}vec\PYZus{}b}\PY{p}{)}
\PY{n}{lena\PYZus{}img\PYZus{}bn}\PY{+w}{ }\PY{o}{=}\PY{+w}{ }\PY{n}{vetor2imagem}\PY{p}{(}\PY{n}{lena\PYZus{}vec\PYZus{}bn}\PY{p}{)}
\end{Verbatim}
\end{tcolorbox}
 
            
\prompt{Out}{outcolor}{ }{}
    
    \begin{center}
    \adjustimage{max size={0.9\linewidth}{0.9\paperheight}}{Trabalho_final_ALC2_files/Trabalho_final_ALC2_100_0.png}
    \end{center}
    { \hspace*{\fill} \\}
    

    \begin{tcolorbox}[breakable, size=fbox, boxrule=1pt, pad at break*=1mm,colback=cellbackground, colframe=cellborder]
\prompt{In}{incolor}{ }{\boxspacing}
\begin{Verbatim}[commandchars=\\\{\}]
\PY{n+nd}{@time}\PY{+w}{ }\PY{n}{α}\PY{p}{,}\PY{+w}{ }\PY{n}{β}\PY{p}{,}\PY{+w}{ }\PY{n}{U}\PY{p}{,}\PY{+w}{ }\PY{n}{V}\PY{p}{,}\PY{+w}{ }\PY{n}{X}\PY{p}{,}\PY{+w}{ }\PY{n}{E}\PY{+w}{ }\PY{o}{=}\PY{+w}{ }\PY{n}{ReconstrucaoLSQR}\PY{p}{(}\PY{n}{A}\PY{p}{,}\PY{n}{lena\PYZus{}vec\PYZus{}bn}\PY{p}{,}\PY{+w}{ }\PY{n}{δ}\PY{p}{,}\PY{+w}{ }\PY{l+m+mi}{10}\PY{p}{,}\PY{+w}{ }\PY{n}{τ}\PY{p}{,}\PY{+w}{ }\PY{n}{ortog}\PY{o}{=}\PY{n+nb}{false}\PY{p}{)}\PY{p}{;}
\end{Verbatim}
\end{tcolorbox}

    \begin{Verbatim}[commandchars=\\\{\}]
\textcolor{ansi-cyan-intense}{\textbf{[ }}\textcolor{ansi-cyan-intense}{\textbf{Info: }}Critério de discrepância não foi
atingido
    \end{Verbatim}

    \begin{Verbatim}[commandchars=\\\{\}]
  0.922215 seconds (395.98 k allocations: 116.481 MiB, 60.98\% compilation time)
    \end{Verbatim}

    \begin{tcolorbox}[breakable, size=fbox, boxrule=1pt, pad at break*=1mm,colback=cellbackground, colframe=cellborder]
\prompt{In}{incolor}{ }{\boxspacing}
\begin{Verbatim}[commandchars=\\\{\}]
\PY{n}{mosaicview}\PY{p}{(}\PY{n}{vetor2imagem}\PY{o}{.}\PY{p}{(}\PY{n}{X}\PY{p}{)}\PY{o}{...}\PY{p}{;}\PY{+w}{ }\PY{n}{nrow}\PY{o}{=}\PY{l+m+mi}{2}\PY{p}{,}\PY{+w}{ }\PY{n}{npad}\PY{o}{=}\PY{l+m+mi}{3}\PY{p}{)}
\PY{c}{\PYZsh{} X1  X3  X5  X7  X9}
\PY{c}{\PYZsh{} X2  X4  X6  X8  X10}
\end{Verbatim}
\end{tcolorbox}
 
            
\prompt{Out}{outcolor}{ }{}
    
    \begin{center}
    \adjustimage{max size={0.9\linewidth}{0.9\paperheight}}{Trabalho_final_ALC2_files/Trabalho_final_ALC2_102_0.png}
    \end{center}
    { \hspace*{\fill} \\}
    

    \hypertarget{comparando-com-um-muxe9todo-burrinho-inguxeanuo}{%
\subsubsection*{\texorpdfstring{Comparando com um método
\sout{burrinho}
ingênuo}{Comparando com um método burrinho ingênuo}}\label{comparando-com-um-muxe9todo-burrinho-inguxeanuo}}

    \begin{tcolorbox}[breakable, size=fbox, boxrule=1pt, pad at break*=1mm,colback=cellbackground, colframe=cellborder]
\prompt{In}{incolor}{ }{\boxspacing}
\begin{Verbatim}[commandchars=\\\{\}]
\PY{n+nd}{@time}\PY{+w}{ }\PY{n}{x\PYZus{}burro}\PY{+w}{ }\PY{o}{=}\PY{+w}{ }\PY{n}{A}\PY{o}{\PYZbs{}}\PY{n}{lena\PYZus{}vec\PYZus{}bn}
\PY{n}{vetor2imagem}\PY{p}{(}\PY{n}{x\PYZus{}burro}\PY{p}{)}
\end{Verbatim}
\end{tcolorbox}

    \begin{Verbatim}[commandchars=\\\{\}]
118.938738 seconds (4.30 M allocations: 4.619 GiB, 0.73\% gc time, 2.83\%
compilation time)
    \end{Verbatim}
 
            
\prompt{Out}{outcolor}{ }{}
    
    \begin{center}
    \adjustimage{max size={0.5\linewidth}{0.5\paperheight}}{Trabalho_final_ALC2_files/Trabalho_final_ALC2_104_1.png}
    \end{center} 
    
    Vemos que de fato, resolver o sistema linear além de ser uma estratégia
ruim em tempo de computação, ainda gera uma solução que é puro ruído.


    % Add a bibliography block to the postdoc
    
\section*{Conclusões}

    Neste trabalho, implementamos o método LSQR para a resolução de
problemas lineares mal postos, aplicando-o na reconstrução de imagens
borradas e ruidosas. Observamos que o método é eficaz na recuperação de
imagens originais, especialmente quando o critério de discrepância é
utilizado para determinar o ponto de parada adequado.

Veja que o método LSQR utilizou poucas iterações, e o tempo computacional foi da ordem de 0.08s, enquanto uma resolução direta do sistema linear levou mais de 118s, uma piora de mais de 1400 vezes no tempo de computação, além de gerar uma imagem que é puro ruído.

A inversão de problemas mal postos é um desafio significativo em muitas áreas da ciência e engenharia. A utilização de uma base sólida da teoria de álgebra linear é essencial para tratar esses problemas de maneira eficaz e robusta, produzindo soluções que tenham significado físico e utilidade prática.

\subsection*{Versão online}

Esse documento está disponível online em \href{https://colab.research.google.com/drive/1yi3uQLi2Yvgg8fH2dJRRlqH82SuueUYv}{Trabalho final - ALC2.ipynb}. 
Esse relatório \LaTeX~foi gerado a partir de um notebook Jupyter usando o pacote \texttt{nbconvert}.
    
    
\end{document}
